\documentclass[]{article}
\setcounter{secnumdepth}{0}
\usepackage[utf8]{inputenc}
\usepackage{authblk}
\usepackage{setspace}
\usepackage[margin=1.25in]{geometry}
\usepackage{graphicx}
\graphicspath{ {./figures/} }
\usepackage{subcaption}
\usepackage{amsmath}
\usepackage{lineno}
\usepackage[export]{adjustbox}
\usepackage{float}
\usepackage[labelfont=bf,justification=raggedright,singlelinecheck=false]{caption}
\usepackage{fancyhdr}
\let\href\undefined
\usepackage[hidelinks]{hyperref}
\usepackage[capitalize]{cleveref}
\usepackage{xspace}
\usepackage{xcolor}
\usepackage{multirow}
\usepackage[normalem]{ulem}
\linenumbers

%%%%%% Running header %%%%%%
\pagestyle{fancy}
\fancyhead[R]{\textbf{machine learning local adaptation}}

%%%%%% Commands %%%%%%
\newcommand{\dfir}{Douglas-fir\xspace}
\newcommand{\ecolocator}{\texttt{EcoLocator}\xspace}
\newcommand{\slim}{\texttt{SLiM}\xspace}

\newcommand{\beginsupplement}{%
    \setcounter{table}{0}
    \renewcommand{\thetable}{S\arabic{table}}%
    \setcounter{figure}{0}
    \renewcommand{\thefigure}{S\arabic{figure}}%
    \renewcommand{\theHtable}{supp.\arabic{table}}%
    \renewcommand{\theHfigure}{supp.\arabic{figure}}%
}

\newcommand{\sbt}[1]{\textcolor{red}{\textbf{SBT: #1}}}


%%%%%% Bibliography %%%%%%
\usepackage[style=authoryear,
citestyle=authoryear,
sorting=nyt,
natbib=true]{biblatex}
\addbibresource{ELrefs.bib}

%%%%%% Title %%%%%%
\title{Supervised machine learning jointly predicts geographic location
and climate of origin from genotypes}

%%%%%% Authors %%%%%%
\author[1*]{Jordan Rodriguez}
\author[2]{Richard Cronn}
\author[1]{Silas Tittes}
\author[1]{Andrew D. Kern}

%%%%%% Affiliations %%%%%%
\affil[1]{Department of Biology, University of Oregon, Eugene, USA.}
\affil[2]{Pacific Northwest Research Station, USDA Forest Service, Corvallis, Oregon, USA.}
\affil[*]{Address correspondence to: jrodrig8@uoregon.edu}

%%%%%% Spacing %%%%%%
\onehalfspacing

\begin{document}

\maketitle

%%%%%% Abstract %%%%%%
\begin{abstract}
Climate change is increasingly disrupting the relationship between 
locally adapted populations and the environments in which they evolved, 
creating an urgent need for tools that connect genomic variation to climate. 
Common-garden and provenance trials remain the gold standard 
for characterizing local adaptation, 
but their time and resource requirements limit how broadly they can be applied. 
Genomic approaches provide a complementary path.
Genotype--environment association (GEA) methods identify environmentally associated loci.
Machine-learning models have also shown that geographic origin can be predicted
directly from genotypes.
Here we introduce \ecolocator, 
a supervised deep neural network that jointly predicts geographic location 
and climate of origin from genotypes. 
Through extensive simulations we demonstrate that
\ecolocator accurately recovers geographic location and environment of origin from genotype data,
and, with SHAP-based feature attribution, identifies adaptive loci more reliably than benchmark GEA methods.
We apply our method to coastal \dfir (\textit{Pseudotsuga menziesii} var. \textit{menziesii}),
where \ecolocator predicts geographic origin ($R^2=0.75$--$0.83$) and climate of origin
($R^2=0.52$--$0.72$) under leave-one-out cross-validation.
Notably, we find that climate predicted directly from genotypes outperforms climate inferred by first
predicting geographic origin, showing that \ecolocator captures genotype--climate signal
that cannot be recovered from geography alone.
Our prediction errors fall within the tolerances used in existing seed-transfer guidelines,
demonstrating that \ecolocator's predictions are ready for practical application, and our approach is
readily extendable to other species and conservation contexts.




\end{abstract}

%%%%%% Main Text %%%%%%

\section{Introduction}

Population genetic theory predicts that the distribution
of genetic variation within and between populations
is shaped by demographic history, spatial structure,
and selective pressures,
and thus carries information about the processes
that generate and maintain it
\citep{lewontin_genetic_1974, hartl_principles_1980, ellegren_determinants_2016}.
In spatially structured, locally adapted populations,
this variation encodes information about both geographic location
and climate of origin.
Geographic structure arises because local mating and limited dispersal
generate isolation by distance,
causing alleles to be more similar among neighboring individuals
than among those farther apart
\citep{wright_isolation_1943, slatkin_isolation_1993}.
Environmental selection, in contrast,
produces allele frequency clines that parallel ecological gradients,
as individuals carrying locally beneficial variants
display higher fitness in particular environments
\citep{endler1977geographic, clausen_experimental_1948, kawecki_conceptual_2004}.
In principle the two signals can be told apart,
since neutral structure follows geographic distance
while adaptive signals follow particular environmental
variables such as temperature, moisture, or seasonality
\citep{rellstab_practical_2015, hoban_finding_2016, tigano_genomics_2016, coop_using_2010}.
In practice, the separation is never clean.
Range expansions, bottlenecks, and drift
can all produce allele frequency gradients that look like adaptive clines
\citep{excoffier_genetic_2009, novembre_spatial_2009}, while spatially 
autocorrelated environmental variables can make real 
adaptive gradients look like isolation by distance, 
inflating false-positive rates in tests
used to distinguish them \citep[e.g.,][]{guillot_dismantling_2013,legendre_should_2015}.
Here we ask whether a single predictive model, trained directly on genomic
data, can learn to tell these two signals apart on its own.

Local adaptation, in which populations evolve higher fitness in their local environments, 
can generate covariation of allele frequencies with environmental gradients 
when spatially heterogeneous selection outweighs the homogenizing effects of gene flow
\citep{savolainen_ecological_2013, lenormand_gene_2002}.
Beyond increasing local fitness,
local adaptation sustains genetic diversity across heterogeneous landscapes
by maintaining quantitative genetic variance and, at some loci, multiple adaptive alleles,
helping populations retain the variation
on which selection can act as environments change \citep{hereford_quantitative_2009}.
The classic models of 
\citet{wright_isolation_1943}, 
\citet{levene_genetic_1953},
and 
\citet{endler1977geographic}
work out how selection, migration, and drift
combine to produce predictable spatial patterns of variation,
and later work has filled in the details.
Even modest selection can create steep genomic clines
when dispersal is limited
\citep{tigano_genomics_2016, yeaman_genetic_2011},
and strong gene flow will swamp adaptive alleles
unless they are tightly linked or of large effect
\citep{slatkin_gene_1975}.

Climate change is now pulling populations out of step
with the environments they are adapted to \citep{lee_ipcc_2023},
and as ranges shift, managers and stakeholders need to know
where populations will persist under conditions
they have not experienced before.
Most of what we have learned about local adaptation comes from
common-garden experiments, reciprocal transplants,
and provenance tests; however these approaches are slow and expensive,
and for most species they have never been done
\citep{aitken_assisted_2013}.
Genomic data are cheaper and faster to collect,
and the field of landscape genomics has grown up around
the effort to connect spatial and environmental variation
with patterns of genetic diversity
\citep{manel_landscape_2003, rellstab_practical_2015}.
In forest trees these methods are already starting to inform
seed transfer and assisted gene flow decisions
\citep{aitken_bemmels_2016}.

Most landscape genetics work relies on genotype--environment association
(GEA) methods, which identify putatively adaptive loci
by testing for statistical associations
between allele frequencies and environmental variables
\citep{rellstab_practical_2015}.
Methods such as \texttt{Bayenv2}
\citep{guenther_bayenv2_2013},
latent factor mixed models
\citep{frichot_lfmm_2013},
and \texttt{BayPass} \citep{gautier_baypass_2015}
test loci one at a time against environmental predictors,
typically while controlling for population structure
through covariance matrices or latent factors.
While widely applied, these methods
are limited in their ability to capture
polygenic adaptation involving many loci of small effect,
and the neutral gradients described above
can produce false positives
\citep{lotterhos_whitlock_2015}.
Newer approaches address some of these limitations
by analysing many loci and environmental predictors simultaneously.
Redundancy analysis (RDA) has proven effective
at detecting multilocus adaptive signatures
\citep{forester_comparing_2018, capblancq_rda_2021},
and gradient forest methods model
nonlinear allele--environment relationships
across the genome \citep{fitzpatrick_gradient_2015}.
More recently, the Weighted-Z Analysis (\texttt{WZA}) introduced
a window-based framework for aggregating
single-locus GEA statistics
into region-level summaries of adaptive signal
\citep{booker_wza_2024}.
All of these methods, old and new, are built to \emph{test}
for associations between genotype and environment.
Spatial population structure enters as a nuisance
to be corrected for, not as something worth modeling in its own right.
So none of them will tell you where an individual came from,
and with the partial exception of gradient-forest and genomic-offset
approaches, which project allele turnover under new climates
\citep{fitzpatrick_gradient_2015}, none of them will predict
an individual's environment of origin from its genotype either.

Machine learning (ML) offers a different way in.
Rather than testing loci one at a time, a supervised model
can use the whole genotype at once to predict
geographic origin or climate of origin directly.
ML methods can also be effective at pulling patterns
out of high-dimensional data with nonlinear interactions among loci
and correlated predictors
\citep{schrider_supervised_2018}, though performance depends 
on how well the training data represent the populations and 
conditions where the model is applied.
\citet{battey_predicting_2020} demonstrated
that a deep neural network, \texttt{Locator}, can predict
the geographic origin of an individual from its genotype alone
by learning from multilocus patterns of isolation by distance.
\texttt{Locator} achieves high spatial accuracy,
but it does not model environmental covariates,
so its predictions give a climate of origin only indirectly, 
through a lookup at the predicted location. While
a location prediction implies a climate of origin once 
paired with a map of historical climate normals, 
that inference assumes the climate at a given location 
stays put — an assumption that erodes as climates shift 
faster than the normals they're being compared against.
Moreover, these prediction offer nothing where future conditions 
have no analog anywhere on today's landscape. 
Predicting climate of origin directly from genotype, 
rather than through an intermediate geographic prediction, 
sidesteps that assumption.

Other recent work brings in ecology more directly.
Gradient forest methods link genomic variation
to environmental predictors to forecast
climate-driven maladaptation \citep{fitzpatrick_experimental_2021}. %ADK -- I replaced these citations. Not sure if these are right. Double check!
In crop genomics, for example, deep learning models that take in weather, soil,
and other covariates improve predictions of how a genotype will perform
across climates
\citep{washburn_predicting_2021, montesinoslopez_review_2021}.
No existing method, though, predicts geography and environment
from genotype together, in one model that can learn from
neutral spatial structure and adaptive environmental associations at the same time.

Here we introduce \ecolocator,
a deep neural network that extends \texttt{Locator}
by jointly predicting geographic location and climate of origin from genomic data.
Trained on georeferenced genotypes paired
with historical climate data,
the model learns multilocus associations
with both geographic position and environmental gradients,
so that it can infer climate of origin directly from genotype
when allele frequencies covary with climate.
We pair the model with a feature attribution method,
SHapley Additive exPlanations
(\texttt{SHAP}) \citep{lundberg_shap_2017},
to see which input SNPs each prediction depends on,
and so which variants contribute most
to adaptive differentiation across the landscape.
We test \ecolocator on simulations to find out
where it works well and where it does not,
then apply it to an empirical dataset of coastal \dfir
(\textit{Pseudotsuga menziesii} var.\ \textit{menziesii}).

Coastal \dfir is a member of the pine family
with a lineage dating to the Paleocene \citep{hermann_genus_1985}
and a natural range spanning extensive climatic
gradients along the North American Pacific coast, from California
to British Columbia. Although most populations are characterized by
low genetic differentiation, high gene flow, and high outcrossing rates \citep{slavov_pollen_2005},
they still exhibit strong signatures of local adaptation,
with adaptive quantitative phenotypes tracking climatic features
such as winter cold temperatures and summer aridity \citep{stclair_genecology_2005}.
A century of provenance trials and seed-transfer research
has shown that trees express maladaptive responses
(e.g., slow growth, poor form, disease susceptibility)
when moved beyond their local source climate,
with home-environment advantage and genotype--environment interactions
structured primarily by winter cold, summer drought,
and summer heat accumulation
\citep{stclair_1912_2020, stclair_genecology_2005, stclair_genetic_2007}.
To minimize maladaptation, seed zones define the geographic limits
within which seedlings can be moved while preserving local adaptation.
\dfir seed zones were formalized for U.S. Forest Service lands
in the 1940s \citep{kummel_forest_1944},
expanded to federal, state, and private lands
in the 1960s and 1970s \citep{usdaforestservice_tree_1973},
and later revised to incorporate climatic factors
and experimental evidence from reciprocal transplant trials
(e.g., \citealp{randall_forest_1996, stclair_seedlot_2022}).
Seed zones also give us something simulations cannot,
which is a standard already in use by managers
for deciding whether a predicted climate is close enough to the true one.
This combination of strong local adaptation and a long history
of ecological and silvicultural study
makes \dfir a natural bridge
between controlled simulations and real-world management applications.


To our knowledge \ecolocator is the first open-source tool
that infers both the geographic origin of a genotype
and the climate it is adapted to,
which is the information managers will need
as climates shift faster than provenance trials can keep up.

\section{\ecolocator}

\ecolocator is an open-source Python package
for predicting the geographic origin and climate of origin
of individual genotypes, accessible via both a command-line 
interface and a programmatic Python API.
\footnote{Source code, installation instructions, and vignettes are available at \url{https://github.com/kr-colab/ecolocator} and \url{VIGNETTE URL}.}
\ecolocator does not depend on fixed assumptions about allele--environment relationships,
but learns directly from high-dimensional genomic
and environmental data (\Cref{fig:architecture}).
By jointly predicting location and environmental variables,
it captures both spatial genetic structure
and signals of local adaptation,
providing a scalable, biologically informed approach
for studying complex evolutionary patterns.
Both the CLI and Python API accept standard genotype formats 
(VCF, Zarr, tab-delimited matrices) as input,
facilitating integration into existing 
genomic workflows whether scripting interactively or from the command line.

\subsection{Preprocessing}

\ecolocator converts genomic data from VCF, Zarr,
and tab-delimited matrix formats
into vectors of minor allele counts per individual
using the scikit-allel \citep{scikit_allel_v1.3.8} and NumPy \citep{numpy_v2.0.0} libraries.
Missing genotypes are imputed by drawing alleles
from a binomial distribution parameterized
by the site's allele frequency,
an approximation that preserves expected variation
while minimizing bias.
Sites can be filtered by a minor allele count threshold
to reduce noise from rare variants
that contribute little predictive signal.
Geographic coordinates are normalized to zero mean and unit variance,
and environmental covariates are standardized similarly
prior to training.

\subsection{Model architecture}

\ecolocator employs a multilayer perceptron with a shared trunk 
and two task-specific output heads, trained under a supervised 
learning framework on labeled datasets of genotypes paired with 
known geographic coordinates and climate variables. Input genotypes 
pass through a batch normalization layer before entering the shared 
trunk, which consists of dense layers with ELU activations and a 
central dropout layer for regularization. The trunk branches into 
two heads --- one predicting 2D geographic coordinates (i.e., latitude and longitude), one predicting 
environmental covariates --- each implemented as a small two-layer subnetwork. 
The network is optimized using the Adam optimizer with a custom multitask 
loss function, combining mean Euclidean distance on the standardized 
latitude and longitude for geographic 
prediction and MSE on the standardized environmental covariates, each scaled by an 
independently specified weight. Concretely, the loss function is:
\begin{equation} \label{eq:1}
    \mathcal{L} = w_{\text{loc}} \left( \frac{1}{n} \sum_{i=1}^{n}
    \sqrt{(\hat{x}_i - x_i)^2 + (\hat{y}_i - y_i)^2} \right)
    + \, w_{\text{env}} \left( \frac{1}{np} \sum_{i=1}^{n} \sum_{j=1}^{p}
    (\hat{z}_{ij} - z_{ij})^2 \right)
\end{equation}
where $n$ is the number of training samples, $p$ is the number of environmental
covariates, $x_i$ and $y_i$ are the true longitude and latitude of sample $i$,
$\hat{x}_i$ and $\hat{y}_i$ are the predicted coordinates, $z_{ij}$ is the true
value of covariate $j$ for sample $i$, $\hat{z}_{ij}$ is the predicted value,
and $w_{\text{loc}}, w_{\text{env}} \geq 0$ are loss
weights (both defaulting to 1.0, see \Cref{fig:loss-weight}).
The network is trained using a standard validation scheme in which a user-specified 
fraction of samples (default 0.9) is used for model fitting, while the remaining fraction is 
held out for validation, enabling convergence monitoring, early stopping, 
and adaptive learning rate reduction to prevent overfitting. 
By explicitly coupling geographic and climatic inference, the model captures 
both neutral spatial structure and environment-associated genomic variation.


\begin{figure}[H]
    \adjustimage{width=1.1\textwidth,center}{figures/EcolocatorARC.pdf}
    \caption{\ecolocator architecture. \textbf{Panel A} shows the training regime,
    with the inputs traversing the network in a supervised fashion; network weights 
    are updated through backpropagation until predictions are optimized toward the
    true value, with the process shown here in blue. Once trained, \ecolocator enters the prediction
    step, shown here in \textbf{Panel B}. Given genotypes alone, the trained network
    jointly predicts climate of origin and geographic location, shown in yellow. }
    \label{fig:architecture}
\end{figure}

In practice, an end-to-end \ecolocator analysis consists of five steps:
\begin{enumerate}
\item Assemble a training dataset consisting of filtered genotype data 
in a supported format (VCF, Zarr, or tab-delimited matrix) 
and a sample metadata file pairing each individual with its 
known geographic coordinates and environmental covariate values. 
Samples to be predicted should have their coordinates and covariates 
left blank.

\item Train the model using \texttt{ecolocator train} (or the Python API), 
specifying the genotype file, sample metadata, and an output directory. 
The model is fit to jointly predict geographic origin and environmental 
covariates, with early stopping and adaptive learning rate reduction applied 
automatically.

\item Evaluate model performance using \texttt{ecolocator loo}, 
which iteratively holds out each sample with known coordinates, trains 
on the remainder, and predicts the held-out individual (leave-one-out cross-validation).
The resulting predictions can be compared against true values to assess spatial and 
environmental accuracy before applying the model to unknown samples.

\item Predict geographic origin and environmental covariates for samples 
with unknown locations using \texttt{ecolocator predict}, providing the 
saved model directory and a genotype file containing the target samples.

\item (\textbf{Optional}) Quantify per-SNP importance for predicted samples 
using \texttt{ecolocator attribute}, which applies SHAP values to 
identify loci driving each prediction.
\end{enumerate}

The predicted coordinates and environmental values from steps 2--4 provide 
continuous estimates of geographic and climatic origin, and SHAP attributions 
from step 5 can be used to identify loci contributing most strongly to each prediction.
More details are given below. 

\section{Methods}

\subsection{Simulating local adaptation in continuous space (non-Wright--Fisher)}

To explore under which evolutionary scenarios \ecolocator performs best,
we used \slim \texttt{v.4.2.2} \citep{haller2023slim4} 
to simulate a spatially explicit non-Wright--Fisher population
experiencing local adaptation of a quantitative trait
that evolves to track a single continuously varying environmental axis.
By varying simulation parameters such as dispersal, genetic architecture of the quantitative trait,
and the strength of stabilizing selection,
the evolutionary trajectory of a given population
can be changed in predictable ways.
For example, when environmental selection is weak
and individuals disperse far distances,
the homogenizing effects of gene flow across the landscape
result in little correlation between genetic variation and environmental variation
\citep{lenormand_gene_2002,kawecki_conceptual_2004,booker_wza_2024}.
The underlying spatial and density-dependent components of the 
simulation were borrowed from \citet{chevy_popgen_ecology_2025},
but are briefly described below for completeness.

\subsection{Fitness and reproduction}

To capture the joint effects of selection, dispersal,
and spatial dynamics,
we simulate individuals in continuous space
on top of an environmental grid (see \Cref{fig:env-maps}).
We initialize each simulation
by randomly placing individuals across the landscape.
Individuals occupy a continuous, square landscape, 
which we divide into an 8\texttimes 8 grid of equal-sized regions, 
each assigned its own environmental value that defines the local phenotypic optimum, 
$\theta$, for individuals within it.
Values are generated using the midpoint-displacement neutral landscape model 
implemented in the NLMpy Python package \citep{etherington_nlmpy_2015},
with a Hurst exponent of zero (H=0),
which removes spatial autocorrelation
and creates a patchwork of optima.
This setup allows us to isolate the effects
of spatial dynamics and dispersal
from those of environment-dependent selection.
Critically, because grid cells have environmental optima that are
uncorrelated at the scale of the grid, 
proximity on the map does not imply 
similarity in environment: 
individuals on opposite sides of the landscape can share 
the same optimum, while individuals in adjacent 
cells can be adapted to entirely different ones.
Geographic and environmental prediction are therefore decoupled;
therefore, any accuracy \ecolocator achieves on the
environmental target cannot arise from geographic
proximity or from neutral spatial structure driven by isolation by distance,
but must instead reflect genotype--environment covariance
generated by local adaptation.

Individuals experience fitness as a combination
of stabilizing selection and density-dependent competition.
Stabilizing selection follows a Gaussian fitness function
centered on the local optimum (\Cref{eq:stabsel}),
where $z$ is the individual's phenotype,
$\theta$ is the optimum in the corresponding grid cell,
and $\sigma_K$ controls the strength of selection:
\begin{equation}
w(z, \theta) = \exp\!\left[-\frac{(z - \theta)^2}{2\sigma_K^2}\right].
\label{eq:stabsel}
\end{equation}
Total fitness is then determined by multiplying this term
with a density-dependent factor $C_i$,
where competition is implemented
as a Beverton--Holt style density regulation function,
\begin{equation}
C_i = \frac{1}{1 + \rho D_i}.
\label{eq:comp}
\end{equation}
where $D_i$ is the local density of neighbors 
around individual $i$, computed as a 
Gaussian-weighted count of nearby individuals, 
and $\rho$ scales the strength of competition. 
Individuals in crowded regions are therefore less 
likely to survive and reproduce.

After survival, individuals disperse
according to a dispersal kernel that 
sets the spatial scale of gene flow,
redistributing individuals across the landscape.
Following dispersal, individuals mate with nearby partners,
so that reproduction depends on both spatial proximity
and local density.
By varying the dispersal parameter
and the standard deviation of the fitness function,
we can impact the strength of local adaptation
in the simulation (\Cref{fig:sim-performance}, \Cref{fig:local-adaptation}).
Together, stabilizing selection, competition, dispersal,
and local mating govern the survival and reproduction of individuals.

Local adaptation can then be quantified
as the mean absolute difference between each individual's phenotype
and its local optimum,
such that smaller values correspond to stronger local adaptation:
\begin{equation}
\mathrm{Local\;Adaptation} \;=\;
\frac{1}{N} \sum_{i=1}^{N} \left| P_i - O_i \right|
\label{eq:localadapt}
\end{equation}
where $P_i$ is the phenotype of individual $i$, $O_i$ is the local optimum at 
individual $i$'s location, and $N$ is the total number of individuals.
Because dispersal, selection strength, and genetic architecture
interact to determine how much local adaptation actually emerges,
this metric lets us measure the realized outcome of each parameter 
combination directly, rather than inferring it from the parameters alone,
so that it can be related to \ecolocator's prediction accuracy. 

\subsection{Genomic architecture and tree sequence recording}
The genetics underlying adaptation are simulated using a 1 Mb chromosome 
carried by each individual, 
along which mutations at quantitative trait loci (QTL) 
arise, with additive effects that determine an individual's phenotype.
Because the simulated chromosome spans only 1 Mb,
a per-base recombination rate of \(10^{-8}\)
produces very little recombination across it
and consequently extensive linkage.
We therefore used \(10^{-6}\),
which yields approximately one expected recombination breakpoint
across the chromosome per meiosis,
reducing long-range LD and allowing clearer evaluation of locus-level
feature attribution (\Cref{fig:ld-decay}).

To explore whether the genomic distribution of adaptive loci
influences the detectability of local adaptation,
we varied the arrangement of QTLs along the chromosome
across three genetic architectures. 
In the first, 
all QTLs were clustered in a single 200 kb window 
(bases 400,000–599,999; spanning 20\% of the chromosome), 
centered at the midpoint of the chromosome. 
In the second, 
QTLs were split across three evenly-spaced 100 kb windows 
(10\% of the chromosome each, 30\% total span), 
centered at bases 250,000, 500,000, and 750,000. 
In the third, 
QTLs were split across five evenly-spaced 100 kb windows 
(10\% of the chromosome each, 50\% total span), 
centered at approximately bases 166,667, 333,333, 500,000, 666,667, and 833,333.

These treatments allowed us to test whether adaptive 
signals are easier to detect when concentrated in 
specific genomic regions versus dispersed more broadly 
(\Cref{fig:sim-performance}, panel a).
We only simulate non-neutral mutations within \slim.
Genomic information is stored in tree sequences
and recapitated using \texttt{pyslim} 1.0.4 \citep{haller2019treeseq}.
Neutral mutations are then overlaid onto the recapitated tree sequences
using \texttt{msprime} \texttt{v.1.1.1} \citep{baumdicker2022efficient},
allowing us to track both adaptive
and neutral background variation separately.

\subsection{Varying dispersal, selection, and genetic architecture parameters}
To evaluate the robustness of \ecolocator
across diverse evolutionary contexts,
we varied three parameters: dispersal,
the strength of stabilizing selection,
and the genomic architecture of adaptive loci (described above).
We picked these three because each one governs
a process that sets how detectable local adaptation is.
Gene flow homogenizes allele frequencies across space,
selection pulls populations in different environments apart,
and genetic architecture decides where in the genome
the adaptive alleles sit.
Dispersal distance (\(\sigma_D\)) sets the spatial scale
of gene flow. Small values keep offspring near their parents
and let neighboring regions differentiate;
large values mix individuals across the landscape.
We tested \(\sigma_D = 0.1, 0.3, 0.5\), which for reference correspond to 2\%, 6\%, and 10\% of the landscape width moved per generation, since the landscape spans $5 \times 5$ spatial units.
The strength of stabilizing selection is set
by the standard deviation (\(\sigma_K\))
of the Gaussian fitness function.
Small \(\sigma_K\) punishes phenotypes far from the local optimum,
while large \(\sigma_K\) tolerates them.
We tested \(\sigma_K = 0.3, 0.5, 0.7\)
for strong, intermediate, and weak selection.
When varying one parameter we held the others fixed:
the dispersal sweep used \(\sigma_K = 0.5\) and genetic architecture 1,
the selection sweep used \(\sigma_D = 0.3\) and genetic architecture 1,
and the architecture sweep used \(\sigma_D = 0.1\) and \(\sigma_K = 0.5\).
Spatial and phenotypic scales in these simulations are defined relative to the 
simulated landscape and environmental gradient rather than to physical distances,
and are therefore best interpreted as ratios.
Mutation and recombination rates, by contrast, are in standard per-base, per-generation units.
Taken together these runs span populations that are highly structured
and strongly adapted at one end,
and weakly structured and weakly adapted at the other,
which is the range over which we wanted to see
where \ecolocator works and where it does not.

\subsection{Simulating local adaptation in discrete space (Wright--Fisher)}

To complement our spatially explicit simulations,
we implemented a second set of forward-in-time simulations
based on the Wright--Fisher framework
described in \citet{booker_wza_2024}.
Simulation design and parameterization closely followed
the implementation provided in the \texttt{WZA} GitHub repository
(\url{https://github.com/TBooker/WZA/tree/master}),
with minor modifications to ensure compatibility
with the current version of \slim.
These simulations were designed to more closely mirror
the evolutionary scenarios used to evaluate \texttt{WZA},
allowing for a more direct comparison
between \ecolocator and existing GEA methods.
Following \citet{booker_wza_2024},
we simulated a two-dimensional stepping-stone metapopulation
composed of discrete demes arranged on a grid,
with a constant population size per deme ($N=100$)
under Wright--Fisher dynamics.
Migration was restricted to neighboring demes at a rate of
$m=0.0375$ per generation in each cardinal direction
(demes at grid edges and corners had correspondingly fewer
migrant sources), producing a pattern of isolation by distance
and spatial genetic structure across the landscape.
Environmental heterogeneity was defined
using the British Columbia (BC) map of climatic variation,
which assigns each deme a local phenotypic optimum
based on discretized environmental values.
In contrast to our main simulations,
where a continuous phenotype is pulled toward a local optimum 
that varies across the landscape and fitness declines with deviation 
from that optimum in either direction,
these Wright--Fisher simulations implement directional selection
at a small number of discrete loci:
at each locus, one allele is simply favored over the other 
depending on the local environment, 
and fitness increases monotonically with the number of copies carried, 
with no penalty for carrying ``too many.''
Specifically, 12 biallelic loci distributed across the genome
experience genic selection,
in which each copy of the derived allele 
multiplies fitness by $(1 + s\theta)$, 
where $\theta$ is the local phenotypic optimum assigned to that
deme and $s$ is a fixed selection coefficient ($s=0.0136$), with
heterozygotes receiving half this effect.
Because $\theta$ is positive in demes where the 
derived allele is advantageous and negative where it is 
disadvantageous, selection at each locus consistently 
favors whichever allele matches the local environment.
This formulation produces strong genotype--environment associations
localized to a relatively small number of genomic regions.

Simulations were run in two phases:
an initial neutral burn-in period
to establish population structure and linkage patterns,
followed by a directional selection phase
during which adaptive alleles increase in frequency
in a spatially heterogeneous manner.
Tree-sequence recording was used to capture
the full genealogical history of each simulation,
and neutral mutations were subsequently overlaid
to generate genome-wide variation.
To prepare as input to \ecolocator, deme coordinates were encoded
as latitude and longitude,
using the row and column indices of each deme in the grid,
and the environmental variable was defined
as the local phenotypic optimum for each deme.
Genotype matrices were extracted from tree sequences
and used as input for both \ecolocator
and downstream GEA analyses.
With few loci under directional selection and structure imposed by discrete demes,
these simulations are a simpler setting than our continuous-space ones,
with genotype--environment associations that are strong
and confined to a small part of the genome.
They give us a useful contrast to the harder spatial scenarios.

\subsection{Feature attribution}

To quantify the contribution of individual genetic features
to \ecolocator's predictions, we use the \texttt{SHAP} Python
package (\texttt{v.$\geq$0.47.0}) \citep{lundberg_shap_2017},
which assigns importance scores
based on each feature's marginal contribution to the model output
across all possible feature coalitions.
Attributions were computed separately for each prediction target
(latitude, longitude, and environmental covariates),
allowing the genetic contributions to be evaluated
conditional on the response variable of interest.
SHAP values were calculated on trained models
using the genotype matrix as input,
attributing predictive signal to individual loci
while accounting for nonlinear interactions among sites.
To complement SHAP-based attribution
and assess the robustness of inferred signals,
we performed a series of permutation-based sensitivity analyses,
testing whether \ecolocator's accuracy depends primarily on the selected QTL window, 
SNPs in linkage disequilibrium with that window, 
or signal distributed across the genotype matrix genome-wide 
(\Cref{fig:permutation}).
In the genotype matrix,
with individuals as rows and SNPs as columns,
we independently shuffled each column's values across individuals prior to prediction,
breaking that SNP's association with each individual's true label
while leaving its allele frequency and every other column unchanged;
because each column was permuted independently, this also destroyed any 
correlation structure among jointly permuted sites.
These permutations included random shuffling
of increasing proportions of sites genome-wide,
targeted permutations of predefined genomic windows
containing selected variants,
and permutations of loci in linkage disequilibrium (LD)
with those windows.
LD partners were identified using pairwise correlation thresholds
(see \Cref{fig:ld-decay}),
and permutations were applied either to the selected window alone
or jointly with subsets of LD-linked sites
stratified by LD strength.
This framework enables systematic evaluation
of how predictive performance responds to disruption
of adaptive loci, correlated hitchhiking regions,
and broader neutral genetic structure,
providing an interpretable link between model predictions,
genetic architecture, and simulated evolutionary processes.
Feature attribution analyses were performed separately
for each simulation regime,
allowing us to assess how the recovery of adaptive loci
depends on the underlying evolutionary scenario.

\paragraph{Empirical FDR scores from SHAP.}
SHAP values do not inherently provide a measure of statistical significance,
as they represent model-derived feature importance scores
rather than hypothesis test statistics.
To enable direct comparison with GEA methods,
which report $p$-values for association,
we transformed SHAP-based rankings into an empirical FDR score framework.
SHAP values were first aggregated across samples
by computing the mean absolute SHAP value per SNP.
Variants were then ranked from highest to lowest
based on this aggregated score.
To enable comparison with GEA-derived $p$-values,
we calculated an empirical FDR score for each variant
using a rank-based framework.
At each rank position $i$,
all variants up to that point were treated as significant,
and the proportion of false positives
(neutral variants)
relative to the total number of variants in the set
was calculated.
This assigns each SNP an empirical FDR score
based on its SHAP threshold.
Because this approach relies on known neutrality labels,
these FDR scores represent an empirical,
supervised estimate of significance
rather than a model-based or permutation-based test.


\subsection{Comparing \ecolocator to other GEA methods}

To benchmark the ability of \ecolocator-derived feature attribution
to identify adaptive loci,
we compared SHAP-based variant rankings
to two genome--environment association (GEA) approaches:
the Weighted-Z Analysis (\texttt{WZA}) \citep{booker_wza_2024}
and the top-candidate method \citep{yeaman_convergent_2016}.
These comparisons were performed across both simulation regimes
described above,
allowing us to evaluate method performance
under distinct evolutionary scenarios.
\texttt{WZA} and the top-candidate method
are both window-based approaches designed to capture
this regional structure in adaptive signal.
\texttt{WZA} combines SNP-level association statistics across genomic windows,
whereas the top-candidate method tests whether a region
contains an excess of highly associated SNPs.
\citet{booker_wza_2024} found that both methods
performed as well as or better than single-SNP methods
(Kendall's $\tau$, \texttt{BayPass}, \texttt{LFMM}, and \texttt{RDA})
under both strong and weak local adaptation,
which is why we use them as our benchmarks here.
Comparing SHAP-based attribution from \ecolocator to these methods
provides a direct test of whether SNP-level importance scores
derived from a trained predictive model
can recover adaptive loci while preserving resolution
at the level of individual variants.

\paragraph{Weighted-Z Analysis (\texttt{WZA}).}
Following \citet{booker_wza_2024},
we first computed per-SNP association statistics
by measuring the correlation
between genotype and environmental values,
and obtained a $p$-value for each SNP.
These SNP-level $p$-values were then transformed
into Z-scores and combined within genomic windows
using a weighted Z-test,
where weights are proportional to expected heterozygosity,
such that SNPs with higher heterozygosity
contribute more strongly to the window-level statistic.
This weighting scheme reflects the greater information content
of polymorphic sites with intermediate allele frequencies.
To account for variation in the number of SNPs per window,
we applied a polynomial-based correction
in which the expected mean and variance
of window-level Z-scores were modeled
as functions of SNP count.
Corrected $p$-values (\texttt{Z\_pVal})
were then obtained under the assumption of normality.

\paragraph{Top-candidate method.}
Following \citet{yeaman_convergent_2016},
we ranked SNPs genome-wide by their association p-values
and designated those in the top 1\% (i.e. above the 99th percentile) 
as candidate SNPs.
For each window, we then used a one-sided binomial test 
to assess whther the observed number of 
candidate SNPs exceeded the number expected by chance, 
given the window's total SNP count and the genome-wide candidate rate. 


\paragraph{Integration across methods.}
For each simulation replicate,
we combined SNP-level SHAP FDR scores,
window-level \texttt{WZA} $p$-values,
and top-candidate $p$-values
into a unified dataset.
Window-based statistics were mapped back to individual SNPs
using SNP-to-window assignments,
allowing direct comparison across methods
at the SNP level.
To obtain a genome-wide assessment,
variant-level results were concatenated
across simulation replicates
within each genetic architecture
and simulation regime.

\paragraph{Performance evaluation.}
Method performance was evaluated
using precision--recall (PR) curves,
which are appropriate for imbalanced datasets
with relatively few adaptive variants.
For each method,
$p$-values (or, for SHAP, empirical FDR scores) were transformed to scores
($-\log_{10}$),
such that larger values correspond
to stronger evidence of association.
Because area under the precision--recall curve (AUC-PR) depends only
on the rank ordering each score induces,
comparing methods this way does not require their underlying
significance measures to share a common statistical interpretation
or null distribution --- only that each score ranks variants
consistently from most to least likely adaptive within its own method.
Precision and recall were computed
across all thresholds, and AUC-PR 
was used as a summary statistic.
The random baseline was defined
as the proportion of adaptive variants
in the dataset. Performance was evaluated separately within each simulation regime,
allowing direct comparison of method behavior
under both spatially explicit and Wright--Fisher evolutionary scenarios.

\subsection{Empirical application to coastal \dfir}

\paragraph{Samples and genotyping.}
Our empirical dataset consisted of open-pollinated seedlings from
348 unrelated coastal \dfir
(\textit{Pseudotsuga menziesii} var.\ \textit{menziesii}) trees
sampled across Oregon and Washington.
The maternal parents of these trees were woodsrun (unimproved) `plus' trees,
identified through field observations as well-adapted to their local environment and exceptional
for diameter, height, and stem form. Leaf-derived DNA from individual trees was genotyped using the 58K \dfir
Axiom Affymetrix SNP array \citep{howe_axiom_2020}.
Genotype calls were generated in Axiom Analysis Suite (v.5.4)
using a high-stringency workflow (dQC $\geq$ 0.90, call rate $\geq$ 0.97),
supplemented by a secondary workflow for samples of slightly lower quality
(dQC > 0.82, call rate $\geq$ 0.90).
We restricted analysis to SNPs classified as PolyHighResolution ---
Axiom's designation for markers with clearly resolved,
high-confidence biallelic genotype clusters --- yielding 18,732
loci and excluding lower-confidence and off-target variants.
Within this marker set, missing genotype rates ranged from approximately 
0.15\% to 9.85\% across samples and were resolved using the 
imputation procedure described above.
The \dfir genome is large and incompletely
assembled ($\sim$16 Gb; \citealp{neale_douglas-fir_2017}),
so genomic positions of array SNPs are unknown,
precluding chromosome- or window-level analyses.
\dfir exhibits low self-fertility and
seeds derived from maternal parents are typically pollinated by multiple nearby
pollen donors. While the pollen pool is heterogeneous, pollen and seed movement
in \dfir is limited (typically $<400$~m),
making the maternal source environment a reasonable proxy for the source of open-pollinated
seedlings.

\paragraph{Climate variables.}
To represent source climates,
we chose 1970--2000 baseline climate normals modeled by
ClimateNA \citep{mahony_global_2022} using source latitude, longitude, and elevation.
We focused on three climate variables shown to be significantly
associated with \dfir growth and local adaptation \citep{stclair_genecology_2005}:
mean cold month temperature (MCMT), which shows the strongest
associations with the timing of summer budset and annual height
increment; and Summer Heat:Moisture Index (SHM) and mean summer
precipitation (MSP; May--September precipitation),
both of which show the strongest associations with
the timing of spring bud break and the magnitude of stem taper.
These variables capture the primary axes of climate variation
across the species' range and correspond directly to the selection
pressures (e.g., winter cold temperatures; aridity) that are
managed and mitigated by USFS seed zones.
To characterize the spatial scale over which each climate variable
varies, we computed a simple spatial correlogram: for each variable,
we standardized true values across all 348 sampled trees, calculated
the product of standardized values for every pair of trees, and
averaged these products within geographic distance bins. We report
the approximate distance at which this average first falls below 0.1
as the effective range of spatial autocorrelation for each variable.

\paragraph{Model training and feature attribution.}
Prior to model training,
we applied the preprocessing steps described above,
and geographic coordinates and environmental covariates
were centered and scaled.
We trained \ecolocator on this dataset
using leave-one-out cross-validation,
enabling direct comparison between
each individual's predicted and known
geographic coordinates and climate values.
We also wanted to know whether the climate predictions
contain information that the geographic predictions do not.
For this we took each predicted location, looked up its climate in ClimateNA,
and compared the accuracy of these indirect estimates
to the climate values \ecolocator predicted directly.
SHAP values were summarized across leave-one-out
folds by calculating the mean absolute attribution
for each of the 18,732 SNPs separately for latitude, longitude,
MCMT, MSP, and SHM. For each prediction target,
we retained the 179 highest-attribution SNPs,
approximately the top 1\% of the ranked SHAP distribution,
as candidate loci for downstream comparison and biological interpretation.
This threshold was used for prioritization rather than as a test of
statistical significance. We then compared pairwise overlap
among the five target-specific SNP sets to determine the extent
to which geographic and climatic predictions relied on shared genomic features.
As a supplemental analysis, we aligned array probe sequences to a recent
draft assembly (DFv5.29) and, restricting to primary,
high-confidence alignments (MAPQ $\geq$ 30), recovered approximate
chromosomal positions for a subset of the genotyped SNPs.
We use these positions only to plot attribution scores along the genome;
they play no role in prediction.

\paragraph{Functional annotation.}
To characterize the potential biological functions represented
among high-attribution loci, SNP identifiers were used to retrieve their source
cDNA sequences. Available sequence annotations included best protein matches
to \textit{Arabidopsis thaliana} and \textit{Populus trichocarpa},
along with Mercator gene descriptions and PANTHER functional classes and family
identifiers. We used these annotations descriptively to identify recurring biological
functions and functional themes among high-attribution loci and to compare these patterns
across prediction targets. Because these functional assignments rely in part on
homology to angiosperm proteins and SHAP quantifies predictive importance
rather than causality, we interpret the resulting annotations as a means of
prioritizing and biologically contextualizing candidate loci rather than as
evidence that individual SNPs are causal targets of local adaptation.

% TODO: Confirm the specific annotation resource/database versions with Rich
% and add citations for PANTHER and Mercator as appropriate.

\paragraph{Validation against USFS seed zones.}
To assess whether prediction errors are biologically meaningful,
we compared \ecolocator-predicted climate values to the climate ranges
within USFS seed transfer zones. In the Pacific west, `first-generation'
seed transfer zones \citep{usdaforestservice_tree_1973} were narrowly
defined to restrict movement and minimize maladaptive responses to winter
cold temperatures. `Second-generation' zones \citep{randall_forest_1996, randall_washington_2002}
were subsequently developed after common garden experiments revealed that
\dfir populations show broader climate tolerance than is reflected in first-generation zones.
We compiled spatial layers for first-generation Oregon and Washington seed zones,
along with climate summaries for seed zones divided into 1000 ft elevation bands,
instead of the original 500-foot elevation bands \citep{shalev_zone_2026}.
This allows us to evaluate results in the context of first-generation seed
zone boundaries (which are widely used in forest management), while allowing for comparisons to climate
transfers that are closer to the experimentally-supported updated second-generation zones.
For each individual tree, we compared \ecolocator climate predictions to known seed zone climates
using two approaches:
(a) a ``fixed-zone'' approach where map-based seed zones were treated as fixed constraints; and
(b) a ``focal point'' approach where the climate range of the local seed zone was used
to define the maximum acceptable climate transfer.
With the fixed-zone approach, a climate prediction is considered a ``match'' if it falls within
the percentile-defined range of the source seed zone and elevation band. This map-based
approach can lead to non-matches when predictions are nearly identical to true values
due to the arbitrary nature of fixed map boundaries.
With the focal point approach, the difference between the predicted
and true climate values is compared to the range of climate variation within the source seed
zone. For example, if a source seed zone--elevation band shows a
MCMT 100th percentile range of 2\textdegree{}C, an \ecolocator MCMT prediction
within 2\textdegree{}C of the true value falls inside the range of allowable transfers
and is considered a ``match''.
In both cases we compared predicted values to the 100th percentile range
(defined by the minimum and maximum values)
and the 98th percentile range (defined by the 1st and 99th percentile values).
The indirect location plus ClimateNA estimates were scored the same way.


\section{Results}

\subsection{\ecolocator prediction accuracy scales proportionally
with strength of local adaptation}

We evaluated \ecolocator
across a range of simulation parameterizations (\Cref{fig:sim-performance}).
Following \citet{booker_wza_2024},
we ran 100 simulation replicates for each parameter combination in our continuous space simulations,
and we ran 30 replicates for the discrete space simulations.
We additionally summarized
the distribution of local adaptation scores
across representative parameter ranges to illustrate
how dispersal, selection strength, and genetic architecture
influence the degree of local adaptation (see \Cref{fig:local-adaptation}).
Under restricted dispersal (\(\sigma_D\) = 0.1) and
stronger selection (\(\sigma_K\)=0.5),
simulated populations became strongly locally adapted
to their environments (LA $<$ 0.1) after 10,000 generations. 
\ecolocator performed best when dispersal was constrained
(\(\sigma_D\) = 0.1, Lon. $R^2$=0.968, Lat. $R^2$=0.967, Env. $R^2$=0.769;
\Cref{fig:sim-performance}, panel a),
while increasing dispersal reduced accuracy
for both geographic and environmental prediction 
(\(\sigma_D\)=0.5, Lon. $R^2$=0.815, Lat. $R^2$=0.814, Env. $R^2$=0.336;
\Cref{fig:sim-performance}, panel a).
This pattern is consistent with increasing gene flow weakening both
isolation-by-distance structure and genotype--environment correspondence.
Weakening selection also reduced environmental prediction accuracy 
\((\sigma_K\)=0.3, Env. $R^2$=0.609;
\(\sigma_K\)=0.5, Env. $R^2$=0.561;
\(\sigma_K\)=0.7, Env. $R^2$=0.421).
In contrast, geographic prediction remained largely unchanged
across selection regimes 
(\(\sigma_K\)=0.3, Lon. $R^2$=0.917, Lat. $R^2$=0.916;
 \(\sigma_K\)=0.5, Lon. $R^2$=0.911, Lat. $R^2$=0.911;
 \(\sigma_K\)=0.7, Lon. $R^2$=0.906, Lat. $R^2$=0.907).
This dissociation is itself informative.
Dispersal, and thus the spatial structure of the population,
was held constant across selection regimes,
and on a landscape with no spatial autocorrelation
in the environment ($H=0$),
geography carries no information about the environmental target.
The environmental signal that \ecolocator recovers here
must therefore come from something other than geographic structure.

Varying genetic architecture produced only a slight difference in local adaptation
(${\Delta}$LA=0.008 between GA1 and GA3),
despite the higher number of QTL mutations arising
under genetic architecture 3 (GA1=25, GA3=35).
Geographic prediction accuracy remained nearly identical across architectures
(GA1 - Lon. $R^2$=0.967, Lat. $R^2$=0.967; 
GA2 - Lon. $R^2$=0.967, Lat. $R^2$=0.967; 
GA3 - Lon. $R^2$=0.968, Lat. $R^2$=0.968),
and environmental covariate prediction accuracy was similarly stable 
(GA1 - Env. $R^2$=0.759; 
GA2 - Env. $R^2$=0.771; 
GA3 - Env. $R^2$=0.787).
Although genetic architecture 1 
exhibited more extreme outliers and genetic architecture 3 showed slightly higher environmental prediction accuracy,
\ecolocator performance was relatively robust 
to the genomic distribution of adaptive loci within the range explored here. 

To test whether this robustness extends to landscapes where environmental 
variation covaries with geography, we ran additional supplemental simulations 
using genetic architecture 3, under two combined dispersal/selection 
regimes: limited dispersal with strong selection 
($\sigma_D=0.1$, $\sigma_K=0.5$), and higher dispersal with 
weaker selection ($\sigma_D=0.5$, $\sigma_K=0.7$) (\Cref{fig:env-map-correlated}). 
Per-replicate leave-one-out cross validation (LOOCV) $R^2$ 
again declined with 
increasing dispersal and weakening selection under this design, 
most sharply for latitude (\Cref{fig:loocv-accuracy}).
\ecolocator also performed strongly under the discrete-space,
Wright--Fisher simulation regime (\Cref{fig:sim-performance}, panel b).
Across 30 simulation replicates,
\ecolocator accurately predicted both deme coordinates
and the environmental optimum assigned to each deme
(Lon. $R^2$=0.986,
Lat. $R^2$=0.985,
Env. $R^2$=0.916).
This is not surprising given how these simulations are built:
migration is restricted to neighboring demes,
so spatial structure is strong,
and local adaptation is driven by a small number of causal loci.
The same features account for the stronger attribution
and GEA results we report for this regime below.
We summarize the full simulation $R^2$ results in \Cref{tab:supptable2}.

\begin{figure}[H] 
    \adjustimage{width=1.1\textwidth,center}{figures/FIG2_simssetups.pdf}
    \caption{Effects of parameter variation on \ecolocator performance. \textbf{Panel a)} 
    shows the continuous-space, non-Wright--Fisher simulation
    parameterizations in the top row, with prediction accuracy for each respective
    parameter shown as boxplots in the bottom row. We varied dispersal
    (\(\sigma_D\) = 0.1, 0.3, 0.5), selection strength
    (\(\sigma_K\) = 0.3, 0.5, 0.7), and genetic architecture
    (GA = 1, 2, 3). Yellow indicates parameter regimes expected to produce
    stronger local adaptation, including limited dispersal and stronger selection,
    whereas blue indicates parameter regimes expected to produce weaker local adaptation, 
    including higher dispersal and weaker selection.
    Boxplots show prediction accuracy, measured as $R^2$ using scikit-learn,
    for longitude, latitude, and environmental targets across the varied
    simulation parameters. Each continuous-space simulation parameter regime
    was replicated 100 times. \textbf{Panel b)} shows the discrete-space,
    Wright--Fisher simulation regime in the top row, with corresponding
    prediction accuracy shown in the bottom row. The Wright--Fisher simulations
    were held under a fixed parameterization and replicated 30 times,
    providing a complementary comparison to the continuous-space parameter
    sweeps. Boxplots summarize \ecolocator prediction accuracy across
    geographic and environmental targets in this discrete-deme framework.} 
    \label{fig:sim-performance}
\end{figure}

\subsection{Feature attribution methods reveal adaptive loci \textit{in silico}}

To evaluate whether feature attribution identifies adaptive loci,
we summarized SHAP values across all simulation replicates under
strong local adaptation ($\sigma_D = 0.1$, $\sigma_K = 0.5$)
and compared attribution scores between non-neutral variants
and neutral variants,
including sites in linkage disequilibrium.
Across all continuous-space genetic architectures,
non-neutral variants consistently received higher SHAP values
than neutral sites,
indicating that \ecolocator preferentially attributes predictive signal
to loci under selection
(\Cref{fig:shap-attribution}, panels a--c).
This separation was evident both in individual simulation examples
and in distributions aggregated across all replicates.

Under genetic architecture 1
(\Cref{fig:shap-attribution}, panel a),
non-neutral variants exhibited significantly higher mean attribution values
($\mu = 0.390$) relative to neutral and LD sites combined
($\mu = 0.292$; $p = 2.21 \times 10^{-4}$).
A similar pattern was observed under genetic architecture 2
(\Cref{fig:shap-attribution}, panel b),
where non-neutral variants again showed elevated SHAP values
($\mu = 0.403$)
compared to neutral and LD variants
($\mu = 0.303$; $p = 5.85 \times 10^{-6}$).
Despite differences in the number of adaptive loci,
genetic architecture 3
(\Cref{fig:shap-attribution}, panel c)
also showed enrichment of SHAP signal at non-neutral variants
($\mu = 0.394$)
relative to neutral and LD sites
($\mu = 0.319$; $p = 8.79 \times 10^{-4}$).

Although effect sizes were modest
(Cohen's $d \approx 0.14$--$0.18$ across architectures),
the consistent shift in the distribution of attribution values
indicates that adaptive loci are preferentially weighted by the model.
Notably,
the magnitude of separation between non-neutral and neutral variants
was similar across genetic architectures,
despite differences in the total number of causal loci
and the extent of linkage disequilibrium.
This observation suggests that feature attribution performance is robust 
to variation in genetic architecture within the range explored here, 
which also held under the spatially correlated environment design 
described above (\Cref{fig:env-map-correlated}): non-neutral loci again showed 
elevated attribution across all three targets, and longitude and 
environment tracked the same genomic regions in a representative 
replicate (\Cref{fig:shap-manhattan-sim}), an expected consequence of the two 
variables following a shared spatial gradient by design.
We observed stronger separation between neutral and non-neutral sites
in the discrete-space Wright--Fisher simulations
(\Cref{fig:shap-attribution}, panel d),
where local adaptation was driven by 12 causal loci.
Non-neutral variants had substantially higher SHAP values
than neutral variants,
with a median attribution value of 0.326
for non-neutral sites compared to 0.132
for neutral sites.
This difference was highly significant
($p = 8.17 \times 10^{-9}$),
indicating that \ecolocator concentrated attribution signal
at causal loci in this discrete-deme simulation regime.
These results show that SHAP-based feature attribution
can highlight loci contributing to local adaptation,
while also capturing signal from linked regions.


\begin{figure}[H]
    \adjustimage{width=1.1\textwidth, max totalheight=0.85\textheight, center}{figures/shapfigMAIN.pdf}
    \caption{Feature attribution across genetic architectures. Each panel on the left shows
    mean absolute SHAP scores across the simulated chromosome for neutral sites 
    (blue), non-neutral sites (pink), and sites in linkage disequilibrium with
    non-neutral variants (yellow). \textbf{Panel a)} represents simulations under
    genetic architecture 1, \textbf{Panel b)} represents simulations under
    genetic architecture 2, and \textbf{Panel c)} represents simulations under
    genetic architecture 3. Across panels a, b, and c, dispersal was held constant
    at (\(\sigma_D\) = 0.1), and the standard deviation of the fitness function was held
    constant at (\(\sigma_K\) = 0.5). 
    % I don't see a description for the horizontal lines, so I am adding this.
    % Please double check.
    In panels a, b, c, horizontal lines in the
    x axis indicate the location of the causal loci.
    \textbf{Panel d)} shows SHAP attributions
    for the Wright--Fisher, discrete-space simulation regime, where 12
    large-effect causal loci (shown as dots on the x-axis) were simulated.
    % There was no description for the boxplots on the right, so I added
    % this. Please check.
    For all simulations, bloxplots on the right show the distribution of
    absolute SHAP values between neutral and non-neutral sites.}

    \label{fig:shap-attribution}
\end{figure}

Permutation-based sensitivity analyses further showed that predictive signal was 
not restricted to the selected QTL window. Permuting the selected window alone 
produced only modest changes in prediction performance, whereas jointly permuting 
the selected window and all LD-linked sites caused accuracy to fall below zero for 
all three targets (X $R^2=-0.114$, Y $R^2=-0.181$, Env. $R^2=-0.267$;
\Cref{fig:permutation}). Because these $R^2$ values are computed out-of-sample on
permuted genotypes, they can be negative: a negative value means the
model's predictions are worse than simply guessing the training-set mean for every
individual, indicating that little or no usable signal remains once those sites are
disrupted. Random genome-wide permutation produced a more
gradual decline in performance, with complete randomization similarly 
eliminating predictive accuracy (X $R^2=-0.507$, Y $R^2=-0.426$, Env. $R^2=-0.425$).


\subsection{Comparison of \ecolocator to other GEA methods}

A notable capability of machine learning approaches like \ecolocator is that,
in this context,
feature attribution methods provide locus-level importance scores
that can be interpreted analogously to genome--environment association statistics.
This allows direct comparison between model-derived attribution scores
and traditional GEA approaches for identifying adaptive loci.
Using SHAP values computed from trained \ecolocator models,
we ranked variants according to their contribution to environmental prediction
and compared their performance to Weighted-Z Analysis (\texttt{WZA})
\citep{booker_wza_2024}
and the top-candidate method \citep{yeaman_convergent_2016}.

To quantify detection accuracy,
we calculated empirical FDR scores
from SHAP rankings, calibrated against known true adaptive status
rather than a null distribution (see Methods),
and evaluated performance using precision--recall curves
(\Cref{fig:gea-comparison}). We refer to this \ecolocator-plus-SHAP pipeline
as \ecolocator\texttimes{}SHAP below, since SHAP attribution is what turns
\ecolocator's trained model into a method that can be benchmarked directly
against GEA approaches like \texttt{WZA}
and top-candidate. The random baseline represents the AUC-PR
expected if variants were ranked at random,
given the proportion of adaptive variants in each simulation regime.
For each simulation regime,
the precision--recall curves were constructed by pooling SNPs
across all simulation replicates,
so the resulting AUC-PR values reflect global discriminatory performance.
We also calculated replicate-level AUC-PR values,
shown as violin plots.
Thus,
the pooled AUC-PR values summarize overall variant-level discrimination
across all replicates,
whereas the violin plots summarize replicate-to-replicate variation
in method performance.

Across all continuous-space genetic architectures,
SHAP-based attribution consistently recovered
a higher proportion of true adaptive variants
than either comparison method in the pooled precision--recall analysis.
Under genetic architecture 1,
SHAP achieved a pooled AUC-PR of 0.1734,
substantially exceeding \texttt{WZA} (0.0911)
and the top-candidate approach (0.0890),
and outperforming the random expectation (0.0426).
Under genetic architecture 2,
SHAP again showed improved performance
(pooled AUC-PR = 0.1870)
relative to \texttt{WZA} (0.0832)
and the top-candidate method (0.0767),
with a random baseline of 0.0514.
Similarly,
under genetic architecture 3,
SHAP yielded a pooled AUC-PR of 0.1786,
compared to 0.0946 for \texttt{WZA}
and 0.0874 for the top-candidate approach
(random = 0.0713).
The replicate-level violin plots showed the same overall ranking.
Under genetic architecture 1, 
\ecolocator\texttimes{}SHAP had the highest mean AUC-PR
(mean AUC-PR = 0.2076),
followed by the top-candidate method
(mean AUC-PR = 0.1281)
and \texttt{WZA}
(mean AUC-PR = 0.1212).
Under genetic architecture 2,
\ecolocator\texttimes{}SHAP again showed the highest mean performance
(mean AUC-PR = 0.2211),
followed by the top-candidate method
(mean AUC-PR = 0.1032)
and \texttt{WZA}
(mean AUC-PR = 0.0995).
Under genetic architecture 3,
\ecolocator\texttimes{}SHAP maintained the highest mean AUC-PR
(mean AUC-PR = 0.2054),
followed by \texttt{WZA}
(mean AUC-PR = 0.1133)
and the top-candidate method
(mean AUC-PR = 0.1049).
Performance differences were most pronounced
at low recall values,
where SHAP maintained substantially higher precision,
indicating that the highest-ranked loci
were enriched for true causal variants.
In contrast,
\texttt{WZA} and the top-candidate method showed rapid declines in precision,
approaching random expectations across much of the recall range.
Notably,
SHAP maintained improved performance
across all genetic architectures,
despite differences in the number
and genomic distribution of adaptive loci.
SHAP-based attribution retained its 
advantage over \texttt{WZA} and the top-candidate method under 
this same spatially correlated design (\Cref{fig:gea-comparison-regimes}).
In the discrete-space Wright--Fisher simulations,
\ecolocator\texttimes{}SHAP showed particularly strong recovery
of adaptive loci,
achieving a pooled AUC-PR of 0.8430
relative to a random baseline of 0.0149
(\Cref{fig:gea-comparison}, column d).
\texttt{WZA} also performed well in this regime
(pooled AUC-PR = 0.4983),
whereas the top-candidate method showed substantially lower recovery
(pooled AUC-PR = 0.0714).
Replicate-level AUC-PR values followed the same ordering,
with \ecolocator\texttimes{}SHAP showing the highest mean performance
(mean AUC-PR = 0.9061),
followed by \texttt{WZA}
(mean AUC-PR = 0.6191)
and the top-candidate method
(mean AUC-PR = 0.1017).

\begin{figure}[H]
    \adjustimage{width=1.1\textwidth,center}{figures/GEAcomparisonfig.pdf}
    \caption{AUC-PR curve comparison between SHAP, \texttt{WZA}, and top-candidate methods.
    Each column represents a different simulation regime. \textbf{Column a)} represents
    genetic architecture 1, \textbf{Column b)} represents genetic architecture 2,
    and \textbf{Column c)} represents genetic architecture 3. \textbf{Column d)}
    represents the Wright--Fisher, discrete-space simulation regime. The top row
    shows violin plots of per-replicate AUC-PR values, where each point represents
    one simulation replicate, the horizontal bar indicates the mean, and the
    annotated value gives the mean AUC-PR for each method; this view captures
    replicate-to-replicate variability in performance. The bottom row shows
    precision-recall curves constructed by pooling SNPs across all simulation
    replicates within each regime, with the annotated AUC-PR giving the area
    under each pooled curve; this reflects global discriminatory performance
    and is not directly comparable to the per-replicate means shown above. In all panels,
    purple indicates \ecolocator\texttimes{}SHAP, yellow indicates \texttt{WZA},
    and blue indicates the top-candidate method.}
    \label{fig:gea-comparison}
\end{figure}


\subsection{Predicting suitable environments
for coastal \dfir using \ecolocator}

Simulations provide a controlled setting to benchmark \ecolocator,
but the ultimate goal is to evaluate its performance in natural systems
where demography, gene flow, selection, and management history
interact in complex ways. To this end,
we applied \ecolocator to 348 coastal \dfir trees
sampled across Oregon and Washington (\Cref{fig:dfir-results}, panel f)
and genotyped at 18,732 SNPs.
For each tree we predicted geographic origin along with three climate
variables (MCMT, SHM, and MSP) known to structure local adaptation in this species,
using leave-one-out cross-validation so that every prediction
could be compared to a known value.

\begin{figure}[H]
    \adjustimage{width=1.1\textwidth,center}{figures/dfirresults.pdf}
    \caption{\dfir case study results. \textbf{Panels a-e.} show true
    vs. predicted scatterplots across all prediction targets. From left to right,
    we show Longitude ($R^2$=0.75), Latitude ($R^2$= 0.83), 
    Mean Cold Month Temperature (MCMT; $R^2$= 0.72), Summer Heat:Moisture Index 
    (SHM; $R^2$= 0.54), and Mean Summer Precipitation (MSP; $R^2$= 0.52).
    \textbf{Panels f-g.}
    both show the sampling range, with \textbf{Panel f.} being the sample locations
    and \textbf{Panel g.} being the seed zone map, with labels only on the seed 
    zones containing samples.
    \textbf{Panels h-j.} show where \ecolocator predictions fall within the 
    established seed zones, with gray boxes showing the environmental covariate's ranges in a
    given zone, '$\times$'s' representing the source sample value, and colored dots representing \ecolocator 
    predictions inside each zone. From left to right, we show MCMT in pink, SHM 
    in yellow, and MSP in orange. 
    }
    \label{fig:dfir-results}
\end{figure}


\subsubsection*{Prediction accuracy}

\ecolocator predicted source latitudes and longitudes with $R^2$
values of 0.83 and 0.75, respectively (\Cref{fig:dfir-results}). Combined latitude
and longitude for each sample gave predicted geographic origins that showed a median error of 51.6km (see \Cref{fig:geo-error}), with
minimum errors as low as 1.2km and maximum errors as high as 646km. Errors
greater than 200km appeared as extreme outliers and were not associated with obvious
geographic (e.g., edge of range; clustered locations) or genetic (e.g., unusual homozygosity)
attributes. Label errors or failed predictions may contribute to these unusually large errors. 

Climate is strongly spatially autocorrelated,
so one might worry that the climate predictions simply ride on the geographic ones.
If that were the case, predicting location and then looking up its
climate in ClimateNA would do at least as well as predicting climate directly.
It does not. Direct estimates from \ecolocator
were more accurate than ClimateNA values at predicted locations
for all three variables (paired Wilcoxon signed-rank test on absolute
error, $n=348$; MCMT $p=1.1\times10^{-6}$, MSP $p=6.0\times10^{-10}$,
SHM $p=1.6\times10^{-13}$; \Cref{tab:supptable1}),
so the model is picking up climate-related signal in the genotypes
that geography alone does not supply. Among the three modeled climate variables,
Mean Cold Month Temperature (MCMT) showed the highest predicted $R^2$ of 0.72, with Mean Summer
Precipitation (MSP) ($R^2$=0.52) and Summer Heat:Moisture Index (SHM) ($R^2$=0.54)
showing lower and comparable prediction accuracies. This pattern is consistent with
differences in the spatial scale of each variable's underlying gradient: across our
sampled trees, MCMT remains spatially autocorrelated out to roughly 150--175~km, whereas
MSP and SHM decorrelate by around 100--125~km, likely reflecting the influence of local
topography (e.g., rain shadows and orographic effects) on precipitation-related variables
\citep{frankson_oregon_2022}. Because \ecolocator's genomic signal is itself spatially
structured, climate variables that vary smoothly across the landscape should be easier
to predict from that structure than variables that vary sharply over short distances.


\subsubsection*{Validation against USFS seed zones}

To assess whether \ecolocator's prediction errors
are biologically meaningful,
we compared \ecolocator-predicted values to the climate ranges
within USFS seed zone--elevation bands,
using the fixed-zone and focal point criteria described in Methods.
Using the fixed-zone approach, we find that 66.1\%
of the MCMT predictions fall within seed zone--elevation band 100th 
percentile ranges, and that the number of matches drops to 49.7\% of predictions when using the 
98th percentile range. The proportion of matches is higher for MSP (86.2\% for the 
100th percentile range) and SHM (87.6\% for the 100th percentile range) (\Cref{tab:table1}). 
These values are all larger than those estimated using ClimateNA predictions from 
\ecolocator geographic predictions (\Cref{tab:table1}). 

\begin{table}[ht]
  \centering
  \caption{Comparison of \ecolocator and Location+ClimateNA predictions across climate variables using the fixed-zone approach.}
  \label{tab:table1}
  \begin{tabular}{lcccc}
  \hline
   & \multicolumn{2}{c}{\ecolocator} & \multicolumn{2}{c}{Location+ClimateNA} \\
  \cline{2-3} \cline{4-5}
  Factor & 100\%ile range & 98\%ile range & 100\%ile range & 98\%ile range \\
  \hline
  MCMT & 66.1\% & 49.7\% & 65.2\% & 48.9\% \\
  MSP  & 86.2\% & 75.9\% & 74.1\% & 57.5\% \\
  SHM  & 87.6\% & 74.7\% & 71.6\% & 55.5\% \\
\hline
\end{tabular}
\end{table}

Using the focal point approach, we find that 93.1\% of the MCMT predictions 
show lower deviations from the true value than the 100th percentile 
range of the source seed zones. The number of matches 
drops to 87.1\% of predictions when using the 98th percentile range 
(\Cref{tab:table2}). Predictions for MSP and SHM showed higher accuracies,
with 98.6\% and 97.7\% of estimates falling within the 100th 
percentile range of the source seed zone, respectively. As was observed with 
predictions within seed zones, \ecolocator showed higher prediction 
accuracy than estimates using ClimateNA predictions from \ecolocator 
geographic predictions (\Cref{tab:table2}).

\begin{table}[ht]
  \centering
  \caption{Comparison of \ecolocator and Location+ClimateNA predictions across climate variables using the focal point approach.}
  \label{tab:table2}
  \begin{tabular}{lcccc}
  \hline
   & \multicolumn{2}{c}{\ecolocator} & \multicolumn{2}{c}{Location+ClimateNA} \\
  \cline{2-3} \cline{4-5}
  Factor & 100\%ile range & 98\%ile range & 100\%ile range & 98\%ile range \\
  \hline
  MCMT & 93.1\% & 87.1\% & 89.9\% & 75.9\% \\
  MSP  & 98.6\% & 93.1\% & 94.5\% & 85.3\% \\
  SHM  & 97.7\% & 92.8\% & 91.9\% & 83.6\% \\
  \hline
  \end{tabular}
\end{table}

Prediction error for different climate variables was not uniform across seed zones,
and patterns differed by MCMT, MSP, and SHM
(\Cref{fig:dfir-results}, panels h-j). In our sample, trees from coastal seed 
zones in Oregon and Washington (e.g., zones 031, 041, 051, 061) showed lower prediction
errors for MCMT than either Cascade or Cascade Crest sources.
Conversely, trees from Oregon and Washington Cascades seed zones showed
lower prediction errors for MSP and SHM than did Coastal or Cascade Crest trees. 
These results demonstrate that \ecolocator
can predict the climate of origin of individual genotypes
at a resolution finer than the climatic breadth of existing seed zones,
the tolerance that actually governs seed transfer decisions.
Our results further illustrate the value of seed zone benchmarking
as a biologically grounded validation framework:
predictions that fall within a seed zone climate envelope
are not just numerically accurate,
but are close enough that forest managers
consider it an acceptable transfer source.

Although the genotype array used for inference carries no genomic coordinates,
our alignment to the DFv5.29 draft assembly placed a subset of SNPs
on chromosomes, and plotting mean $\lvert\mathrm{SHAP}\rvert$ against these positions
lets us view attribution scores in genomic context for each of the
five prediction targets (\Cref{fig:shap-manhattan-dfir}).
Visually, several genomic regions of \Cref{fig:shap-manhattan-dfir} show elevated attribution for both longitude and
MCMT simultaneously, consistent with the fact that Mean Cold Month
Temperature covaries with longitude across coastal \dfir's range,
with coastal populations experiencing milder winters than those toward 
the Cascade crest---the same pattern we recover in simulation when 
environment is spatially correlated with geography (\Cref{fig:env-map-correlated}, \Cref{fig:shap-manhattan-sim}). 

% I like the information here a lot! It is very nice to see these relationships!
We next asked whether the same genomic features contributed strongly across geographic 
and climatic prediction targets. Comparing the 179 highest-attribution SNPs for 
each target revealed substantial variation in pairwise overlap (\Cref{tab:dfir-shap-overlap}). 
Only 13 SNPs were shared between latitude and MCMT (7.3\%), 
whereas longitude and MCMT shared 97 SNPs (54.2\%). 
The strongest overlap occurred between MSP and SHM, which shared 133 SNPs (74.3\%). 
These overlap patterns mirror relationships among the prediction targets themselves. 
MSP and SHM both capture variation related to summer moisture availability, a 
relationship that is reflected by the largest overlap among high-attribution SNPs. 
Similarly, MCMT varies strongly along the longitudinal transition from relatively 
mild maritime environments toward colder environments closer to the Cascade Range 
and shares substantial genomic signal with longitude. In contrast, 
latitude and MCMT shared relatively few high-attribution loci (7.3\%), 
consistent with latitude representing a north--south geographic axis rather
than the longitudinal winter-temperature gradient captured by MCMT. Because
MSP and SHM are themselves strongly correlated (SHM is calculated in part
from MSP) and MCMT covaries with longitude across the sampled range, some of
this overlap likely reflects intercorrelation among the targets themselves
rather than independent biological signal specific to each; the overlap
patterns are informative about which loci the model draws on for correlated
targets, not evidence that those loci are independently causal for both.

Functional annotation of the high-attribution SNPs revealed additional
differences among prediction targets. Because Axiom array probes are designed
from cDNA sequence (see Methods), each genotyped SNP already has a known
transcript-level origin, so ``loci'' here refers to SNPs anchored within an
expressed gene's transcript sequence, rather than an anonymous genomic
position matched to the nearest annotated gene. Among loci with informative annotations,
latitude-associated candidates included putative homologs involved in 
phytochrome-mediated light signaling, circadian and seasonal regulation, 
and developmental timing, including \textit{PIF3}, \textit{FPA}, WNK-family kinases, 
and several MYB transcription factors. MCMT- and longitude-associated candidates 
included putative \textit{ELF3} and DREB/AP2-family transcription factors, 
SPL-family regulators, molecular chaperones including DnaJ proteins, 
membrane lipid and desaturase-associated genes, and genes involved in 
carbohydrate and trehalose metabolism. These annotations represented 
functions associated with seasonal regulation, temperature-stress response, 
membrane remodeling, and acclimation \citep{lee_photoperiodic_2012, chen_acyl_2013}. 

A distinct but internally consistent set of functions was observed 
among the strongly overlapping MSP and SHM candidate sets. 
These included putative PYL-family ABA receptors, PP2C phosphatases \citep{park_abscisic_2009}, LEA proteins, 
an osmosis-responsive factor, GST-family proteins, membrane transporters, 
and other genes associated with cellular stress protection. 
Additional candidates across the MSP and SHM sets were associated with heat-shock, mitochondrial stress, and redox-response functions.
Although many high-attribution loci were 
poorly annotated or had functions that were difficult to relate directly to climate, 
the interpretable candidates repeatedly represented processes associated with 
environmental timing, temperature acclimation, water-stress signaling, 
membrane remodeling, and oxidative-stress response. % TODO: Add final Supplementary Data S1 file and verify all named candidate annotations against the complete annotation list.
The complete set of high-attribution SNPs and associated functional annotations is provided in Supplementary Data~S1.


\section{Discussion}

\subsection{Joint prediction of geography and climate reveals
spatial structure and local adaptation}

It is well established that multilocus genomic variation
encodes both neutral spatial structure
and signals of environmental adaptation.
Genetic principal components mirror geography
\citep{novembre_genes_2008},
allele frequencies covary with environmental gradients
under local selection
\citep{coop_using_2010, savolainen_ecological_2013},
and isolation by environment can generate genetic structure
independent of geographic distance
\citep{wang_isolation_2014}.
Continuous spatial structure leaves particularly strong
and pervasive signatures in genomic data,
influencing diversity statistics, demographic inference,
and genotype--environment associations
in ways that discrete-deme models fail to capture
% The Battey2020 citaions is rendering weirdly to me, with the full
% names of Peter and Andy. Not a big deal, just FYI.
\citep{battey_space_2020, chevy_popgen_ecology_2025}.
The method \texttt{Locator} \citep{battey_predicting_2020} demonstrated
that these spatial signatures are rich enough
to predict the geographic origin of individuals
from genotypes alone using deep neural networks,
achieving high accuracy across diverse taxa.
\ecolocator extends this framework
by jointly predicting geographic location
and climate of origin,
unifying spatial and environmental inference
in a single supervised model.
This allows us to ask not only where an individual comes from,
but what environmental conditions it is adapted to,
and how the recovery of each signal
depends on evolutionary parameters.
In our simulations, increasing dispersal hurt both geographic and environmental
prediction, but weakening selection hurt only environmental prediction.
This decoupling is consistent with the expectation
that these two signals have distinct evolutionary origins:
geographic prediction depends on isolation-by-distance patterns
driven by limited gene flow,
whereas environmental prediction requires
allele frequency clines that track ecological gradients---a
signal that is eroded
when gene flow overwhelms local selection.
Prediction accuracy barely changed across the three genetic architectures
we tested, though we note that these architectures cover only a small part
of the genetic complexity expected in natural populations.

Our primary continuous-space simulations deliberately decoupled
environment from geography by imposing zero spatial autocorrelation on the
environmental layer ($H=0$), a much sharper decoupling than real landscapes
typically show: in our \dfir dataset, climate variables remain spatially
autocorrelated out to roughly 100--175~km (see above), whereas environmental
gradients commonly covary with geographic structure in natural landscapes
more generally \citep{meirmans_ibd_2012}.
We therefore repeated our analysis under a spatially correlated
environmental gradient to ask whether the same general patterns
persist when geography and environment share spatial structure.
Prediction accuracy again declined under higher dispersal
and weaker selection, indicating that the relationships observed
in the decoupled simulations persist under this more realistic
spatial configuration.
However, when geography and environment covary,
shared predictive signal cannot necessarily be attributed uniquely
to demographic structure or local adaptation,
an important limitation that we return to below.

The Wright--Fisher results show what \ecolocator can do
when the problem is easy. Discrete demes, migration restricted to
nearest neighbors, and directional selection on a dozen causal loci
give the model a clearly structured task, and relative to the size of
the simulated grid, migration is slow enough that spatial and
genotype--environment structure persist. In our continuous-space
simulations, by contrast, diffuse genetic architectures, high
dispersal, or weak selection erode these signals; because dispersal in these
simulations is a non-trivial fraction of the landscape width per
generation (2--10\%, see Methods), individuals mix across the range
more readily than in the Wright--Fisher grid. It is likely this relative scale of
dispersal to range size, rather than continuous space \textit{per se},
that drives the difference, though we have not tested whether a
proportionally larger continuous-space landscape would recover
Wright--Fisher-like performance. The contrast between the two regimes
is nonetheless a reminder that the difficulty of the prediction
problem depends on the underlying biology \citep{booker_wza_2024, dauphin_rethinking_2023}.

The \dfir case study extends these findings to a natural 
system where the ground truth is not fully known. 
\ecolocator's predictions of geographic origin and 
climate variables in coastal \dfir are consistent with the 
species' well-documented patterns of local adaptation along steep 
latitudinal and elevational gradients in western 
North America \citep{aitken_assisted_2013}. 
That the model performs well in this system
despite low population differentiation \citep{krutovsky_estimation_2009}, 
incomplete genomic resources \citep{neale_douglas-fir_2017}, 
and a relatively modest number of SNPs, 
suggests that the joint prediction framework can recover 
informative genomic signal even under challenging empirical conditions. 
Feature attribution provides additional resolution into how this 
predictive signal is distributed across geographic and environmental
targets, which we discuss in detail below.


\subsection{Feature attribution as a tool for identifying candidate adaptive loci}

SHAP-based feature attribution provides a window
into which genomic regions drive \ecolocator's predictions. That
\ecolocator's direct climate predictions outperformed a
geography-then-ClimateNA lookup in \dfir (see above) already indicates
that this signal reflects genuine genotype--environment covariance
rather than geography alone; feature attribution lets us ask which
loci carry that signal. Our simulation results demonstrate
that non-neutral variants consistently receive
higher attribution scores than neutral sites,
confirming that the model preferentially weights loci
under selection.
This signal persists across genetic architectures
despite modest effect sizes,
indicating that SHAP can detect the distributed signature
of polygenic adaptation
even when individual locus effects are small.

However,
the modest effect sizes also highlight
an inherent challenge:
under polygenic adaptation,
the signal at any single locus is subtle,
and attribution scores for adaptive and neutral loci
overlap substantially.
This is consistent with the broader finding
that GEA methods are generally underpowered
under weak selection,
and that detecting most causal loci
requires tolerating a substantial false positive rate
\citep{lotterhos_whitlock_2015, booker_wza_2024}.
\ecolocator's feature attribution faces the same trade-off:
thresholding on SHAP values to identify candidate loci
will inevitably involve a balance
between sensitivity and specificity
that depends on the unknown strength of selection
and the genetic architecture of adaptation.

The continuous-space simulations therefore represent
a comparatively diffuse attribution problem,
with adaptive effects distributed across loci
and additional predictive signal extending into linked sites
through hitchhiking.
Linkage disequilibrium further complicates interpretation,
as SHAP values are elevated not only at causal variants
but also at linked neutral sites.
This is expected,
because LD spreads adaptive signal
across genomic neighborhoods.
The recombination-rate comparison illustrates this effect directly:
LD extended much more broadly beyond the selected QTL window
under \(10^{-8}\) than under the \(10^{-6}\) rate
used for the analyses reported here
(\Cref{fig:ld-decay}).
Together, distributed adaptive effects and signal shared among linked sites
help explain both the modest separation between neutral
and non-neutral SHAP distributions
and the difficulty of resolving individual causal sites.

The permutation-based sensitivity analyses
complement SHAP by providing a different lens
on the same question.
By disrupting specific genomic regions
and measuring the resulting loss in prediction accuracy,
these analyses show how much predictive information is carried
by the selected window,
by linked sites,
and by genome-wide background variation.
Prediction declined much more strongly
when the selected window and LD-linked sites were permuted together
than when the selected window alone was disrupted
(\Cref{fig:permutation}).
These findings suggest that \ecolocator is not relying
on the selected window alone,
but on a broader LD-associated haplotypic signal
that can obscure the distinction between causal and linked variants.

Taken together, these results suggest that SHAP is most reliable for
identifying candidate genomic regions or haplotypic blocks associated
with local adaptation, rather than for pinpointing individual causal
mutations. This distinction matters even more in empirical systems,
where LD can reflect physical linkage, demographic history, selection,
population structure, or unobserved ancestry, and where the true
causal variants are rarely known.

Notably,
SHAP-based attribution from \ecolocator
shows greater power to identify adaptive loci
than \texttt{WZA} \citep{booker_wza_2024}
and top-candidate approaches
in our simulation comparisons
(\Cref{fig:gea-comparison}).
This advantage likely reflects the fact
that SHAP attribution operates on a model
that has already learned the joint relationship
between genotypes,
geography,
and environment,
rather than testing each locus independently
against a single environmental predictor.
By conditioning on a trained multitask model,
SHAP can capture contributions from loci
whose individual marginal correlations with environment are weak
but whose combined effect is detectable
in the context of the full genotype.
This shows that supervised prediction
followed by \textit{post hoc} feature attribution
can be a powerful complement to traditional GEA approaches,
particularly under polygenic adaptation
where per-locus signals are subtle.
At the same time, 
the stronger performance of both SHAP and \texttt{WZA}
in the Wright--Fisher simulations suggests that 
locus-identification performance can improve 
when adaptive signal is strong and localized,
whereas more diffuse adaptive architectures 
present a more difficult inference problem.

One caveat becomes important when feature attribution moves
from simulations to real landscapes.
 Our continuous-space simulations use an environmental layer with no spatial autocorrelation ($H=0$), so that environment and geographic structure share as little signal as possible.
In nature the two usually covary, and when they do,
attribution alone cannot tell us whether a high-attribution locus
reflects environmental selection, demographic structure,
linkage to a selected locus, or some mix of these
(we return to this under \hyperref[sec:strengths-limitations]{Strengths and limitations})).
High-attribution loci should therefore be read as candidate
genomic regions associated with the predicted gradient,
not as confirmed targets of adaptation.

The \dfir results nevertheless tell us something
about how this signal is structured.
The high-attribution SNP sets overlapped most for pairs of targets
that track related gradients, MSP with SHM and longitude with MCMT,
and hardly at all for latitude and MCMT (\Cref{tab:dfir-shap-overlap}).
As noted above, some of this overlap likely reflects correlation among
the target variables themselves rather than independently causal shared
loci. Still, the low overlap between the more weakly correlated targets
(latitude and MCMT, 7.3\%) suggests that \ecolocator is not simply
handing the same spatially structured loci high attribution for every
target, and that a meaningful part of the signal is specific to each
axis of geographic or environmental variation.
The functional annotations echo this pattern. Rather than all five targets
converging on the same set of candidate genes, each target's high-attribution
candidates are dominated by a distinct functional theme, and these themes
line up with what is already known in coastal \dfir. Latitude candidates
were dominated by genes involved in light sensing and seasonal
developmental regulation, which fits the known effects of photoperiod
and temperature on growth cessation \citep{ford_photoperiod_2017}. Longitude
and MCMT candidates were dominated by genes involved in seasonal regulation
and temperature acclimation, in line with genomic evidence for polygenic
cold adaptation in this species \citep{delatorre_polygenic_2021}.
MSP and SHM candidates repeatedly involved ABA-mediated water-stress signaling,
osmotic protection, and redox response, which matches common-garden evidence
for genetic variation in drought and cold hardiness along
summer-precipitation and winter-temperature gradients
\citep{bansal_tolerance_2016}. Viewed together,
these patterns suggest that environmental timing, including the timing of growth, 
dormancy, acclimation, and stress response, may be an important component of the 
genotype--environment signal captured by \ecolocator in \dfir.

\subsection{Implications for \dfir management and forestry}

For natural resource management,
\ecolocator addresses a practical need:
identifying the geographic and climatic source of biological materials
under historical, current, or projected future climates.
Seed transfer guidelines have been used on U.S. Forest Service lands
for over a half century \citep{kummel_forest_1944}, and comparable guidelines 
extend back nearly a century in Europe \citep{myking_historic_2016}. 
Despite the history and ubiquity of seed transfer zones, 
large-scale restoration efforts with non-local materials have been a 
common occurrence in forestry, both in the management of native forests
and in introduced species trials.
These unintended `local adaptation experiments' provide clear examples of 
maladaptive responses, but they are usually associated with a poor
understanding of the seed source. Tools like \ecolocator can
provide estimates of the source location and environment,
making it possible to interpret maladaptive responses in large-scale 
plantings in light of geographic or climatic transfer distances. 
Applying \ecolocator to trials of suspected `off-site' origin could
substantially expand our knowledge of growth and
fitness responses of large climate transfers, as designed provenance tests with 
`off-site' materials are limited in number for most tree species.

Similarly,
forest tree provenance tests established outside native ranges
typically have the goal of identifying high-performing genotypes and geographic
sources where additional material can be selected.
Incomplete records or missing tree identities limit the utility of these tests
by making it difficult to match similar climates to identify pre-adapted seed for screening
and breeding.
Efforts to combat illegal logging now increasingly rely
on DNA-based methods to identify the geographic source of tissues,
such as wood.
Adding climatic source information to these investigations
can help investigators discriminate among geographically similar
but climatically distinct sources (such as those separated by elevation).
Finally,
traditional reforestation programs based on seed transfer guidelines
and common-garden trials are increasingly strained
by the pace of climate change \citep{aitken_bemmels_2016},
and new tools are needed to manage climate uncertainty.
Genomic tools that predict the climate of origin of a given genotype
offer a unique strategy for answering these questions,
enabling managers to evaluate whether replanting resources
are well-matched to the conditions they will encounter
at a given site.

Coastal \dfir is a natural test case for this application.
Its economic importance,
extensive provenance trial data,
and strong signals of local adaptation along climate gradients
provide both the motivation and the validation framework
for genomic prediction tools.
By predicting climate of origin alongside geographic location,
\ecolocator provides information that is directly relevant
to seed transfer decisions:
two seed sources from similar geographic locations may differ
in their predicted climate of origin,
reflecting adaptive differences that might be missed
by geography-based guidelines alone.
By comparing \ecolocator predictions to \dfir seed zones,
these predictions become interpretable in definable management terms.
For example,
the most recently proposed seed zone--elevation bands
in Oregon and Washington span approximately
2.5\textdegree C in MCMT across the 98th percentile range \citep{shalev_zone_2026}.
Our results suggest that \ecolocator predictions fall
within these climate tolerances 87.1\% of the time with \dfir. 
Importantly,
this level of accuracy was achieved with a modestly sized dataset
of 348 trees and approximately 18,000 SNPs,
and without any use of genomic position.
\ecolocator operates on the genotype matrix alone,
so neither a contiguous assembly nor an ordered marker map is required---an
advantage in species whose genomic resources remain incomplete.
Although additional sampling is not evaluated empirically here, downsampling in
our simulations (see \Cref{fig:downsampling}) and
previous work with Locator suggest that geographic prediction accuracy
increases with increased sampling density and,
to a lesser extent,
increased genomic sampling \citep{battey_predicting_2020}.
Similar improvements should be possible for \ecolocator,
particularly in regions where current predictions are limited
by sparse sampling or rapidly changing environmental gradients.

At the same time,
the interpretation of \ecolocator predictions depends on
the sampling history and biological structure of the material being analysed.
The \dfir individuals used here were designated as ``plus'' trees,
meaning they were identified through provenance trials
or field observations as being particularly well-adapted
to their local environments.
This selection criterion may strengthen genotype--environment correspondence
and aid climate prediction,
but it could also overestimate performance relative to
less-characterized or more heterogeneous samples.
The spatial distribution of prediction error also highlights
important considerations for applying \ecolocator in complex landscapes.
For example, the highest prediction errors for MCMT in our dataset
were concentrated near the eastern crest of the Cascade Range,
where several non-mutually exclusive factors may contribute to reduced prediction accuracy.
Notably, within this crest-adjacent zone, error was significantly higher
in the southern Cascades than in the northern Cascades
(mean $\lvert\text{error}\rvert$ 1.57$^\circ$C vs.\ 1.03$^\circ$C;
Mann--Whitney $p=2.7\times10^{-4}$), the opposite of what would be expected
if steep, complex terrain alone were driving these errors, since the
northern Cascades are topographically steeper \citep{haugerud2009}. This pattern instead
points to introgression with interior \dfir
(\textit{Pseudotsuga menziesii} var.\ \textit{glauca}) \citep{stclair_genecology_2005}
as the more likely explanation: coastal-interior admixture is more
prevalent toward the southern Oregon Cascades, near the Klamath--Siskiyou
transition zone, and trees with greater interior ancestry may show higher
prediction error when evaluated with a model trained primarily on coastal
\dfir genomic and climatic data. Steep elevational gradients and low local
sampling density near the crest, and positional inaccuracies in older
Township--Range--Section-derived seed source records, may also contribute
to reduced prediction accuracy in this region.
These patterns suggest that prediction errors
point to biologically meaningful features of the study system,
including ancestry transitions,
complex topography,
and uncertainty in historical source records \citep{rehmann2024evaluating}.

\subsection{Strengths and limitations}
\label{sec:strengths-limitations}

A key strength of \ecolocator is its multitask design, 
which simultaneously models geographic and environmental variation. 
This joint framework allows the model to learn both shared and target-specific 
components of genomic variation and, when sufficient information is present in the 
genotype data, to distinguish signals associated with geographic structure from those 
associated with environmental variation. This design also bears on a well-documented
challenge of conventional GEA methods. When population structure aligns with
environmental gradients, as it often does in nature, methods that correct for
structure can inadvertently remove true adaptive signal along with neutral
confounders \citep{lotterhos_recombination_2019, booker_wza_2024}.
\ecolocator takes a different route: rather than regressing geographic structure
out, it models geographic and environmental variation as related but distinct
prediction tasks. However, this does not resolve the confounding. When climate
covaries with geography, non-causal but spatially structured loci can still
receive high attribution for climate simply because they predict location, and
\ecolocator does not attempt to ask whether a given locus carries environmental
signal beyond what spatial structure explains, as more classic GEA methods do.
Attribution from \ecolocator is best interpreted as identifying loci that are
informative for the joint prediction task, with the question of which of those
are adaptive left to the follow-up analyses discussed above.

Furthermore, the multitask framework
is naturally suited to systems
where adaptation occurs along multiple environmental axes.
\citet{booker_wza_2024} noted that window-based GEA methods
can lose power when adaptation involves
semi-independent environmental gradients,
because the LD patterns around causal alleles
may differ across gradients.
\ecolocator's ability to predict multiple environmental covariates
simultaneously may help retain signal in such cases,
though this remains to be tested explicitly.
Within the range of architectures tested,
\ecolocator performed comparably across simulations
with clustered and distributed QTL arrangements,
and it requires no prior specification of candidate loci
or assumptions about the genetic basis of adaptation.

However, several limitations should be noted.
First, \ecolocator is a supervised method
that requires georeferenced samples paired
with environmental data for training.
Its predictions are therefore bounded
by the geographic and environmental range of the training set,
and extrapolation to novel environments
or unsampled regions should be interpreted cautiously.
Second, like all machine learning approaches,
the model learns correlations rather than causal mechanisms.
High prediction accuracy for an environmental variable
does not by itself demonstrate that selection
on that variable has shaped genetic variation.
As \citet{booker_wza_2024} emphasize,
GEA signals do not constitute a test of local adaptation:
correlated environmental predictors,
spatially structured neutral variation,
allelic surfing during historical range expansions,
and adaptation to unmeasured but environmentally correlated biotic
factors, such as local mycorrhizal or soil microbial communities,
can all generate allele frequency patterns
that mimic adaptive signals
\citep{klopfstein_surfing_2006, novembre_spatial_2009}.
Isolation by distance alone can inflate
genotype--environment correlations
when environmental gradients are spatially autocorrelated
\citep{meirmans_ibd_2012}.
\ecolocator's joint prediction framework
does not resolve this fundamental ambiguity;
high environmental prediction accuracy
could reflect true adaptive signal, demographic history,
or both.
Third, our simulations use a simplified landscape
with a single environmental axis
and a modest genome size.
Natural systems involve many correlated environmental variables,
complex demographic histories,
and genomes orders of magnitude larger,
all of which may affect model performance in ways
not fully captured here.

\subsection{Future directions}

Several extensions of this work merit exploration.
First, our simulations assume a purely additive genetic architecture,
and whether prediction accuracy and attribution hold up under
non-additive genotype--phenotype maps remains untested.

Second, expanding the environmental complexity
of simulations to include multiple correlated environmental axes,
more complex spatial covariance among those axes,
and temporally varying selection would
better approximate the conditions
under which \ecolocator will be applied in practice. 
A particularly promising application is using 
differences in environmental prediction accuracy
to identify which environmental gradients
are most strongly reflected in genomic variation.
Many landscape genomic and genomic-offset approaches
consider multiple environmental predictors,
but not all environmental axes are expected
to contribute equally to local adaptation.
Because \ecolocator predicts each environmental target separately,
variation in prediction accuracy among targets
provides information about the strength
of genotype--environment correspondence along each axis.
For species with limited ecological or common-garden information,
this could help researchers and managers prioritize
which environmental variables warrant the greatest consideration
when designing management strategies,
seed-transfer guidelines,
or subsequent experimental studies. The same logic extends to geographic
prediction: poor geographic accuracy can indicate weak population
structure, high gene flow, or limited spatial signal in the sampled
genomic data, making prediction accuracy itself a diagnostic signal
for species that lack common-garden data or prior characterization.

Third, the framework could also be extended
to future climate applications.
\ecolocator-estimated climates of origin could be compared
with projected future conditions
to evaluate source--site climate mismatch and 
identify candidate sources 
for assisted migration or restoration.
This would connect genotype-based inference
of historical climate associations
with decisions about where biological material
may be best suited under changing conditions. 

Finally, systematic comparison
with a broader suite of GEA methods
across a wider range of evolutionary scenarios
would help delineate the conditions
under which \ecolocator's joint prediction approach
offers the greatest advantage
over single-task or univariate alternatives.

\section{Acknowledgments}
We thank Peter Ralph and \'Angel Rivera-Col\'on for their helpful comments on the manuscript,
Amanda de la Torre and her lab for providing the genomic positions for the \dfir{} data,
and all members of the Kern-Ralph colab who provided feedback and support throughout this project.
This work was supported by NIH grants R35GM148253 and R01HG010774 to ADK.


\printbibliography

\newpage
\beginsupplement
\section*{Supplementary Materials}

%%%%%% Supplemental Figure 1 %%%%%%
\begin{figure}[H]
    \centering
    \begin{subfigure}[t]{0.49\textwidth}
        \centering
        \includegraphics[width=\textwidth]{figures/suppfig1_lossweight_env.png}
        \caption{}
        \label{fig:loss-weight-env}
    \end{subfigure}
    \hfill
    \begin{subfigure}[t]{0.49\textwidth}
        \centering
        \includegraphics[width=\textwidth]{figures/suppfig1_lossweight_loc.png}
        \caption{}
        \label{fig:loss-weight-loc}
    \end{subfigure}
    \caption{\textbf{Loss weight cross-validation.}
    To determine the default loss weight value for location and
    environmental prediction, we tested sliding this value at the command
    line for each output task (location prediction and environmental
    covariate prediction) on the empirical \dfir{} data, while holding
    the other value at 1. In panel \textbf{a)}, we see that when the loss weight is
    held at a value of 0 for environmental covariate prediction, \ecolocator
    essentially ``focuses'' entirely on the location prediction task,
    reflected in the low $R^2$ scores for all three environmental targets and
    high prediction scores for location. We see the same effect in the
    opposite direction in panel \textbf{b)}, where we held the loss weight of the
    environmental output head at 1 and scaled the loss weight for location
    prediction from 0--1. Again we see that when loss weight for location
    prediction is held at 0, the model prioritizes predicting on
    environmental covariates. Reassuringly, this analysis also demonstrates
    that \ecolocator recovers the true values for environmental covariates
    without over-reliance on the spatial information housed in the location
    values. Lastly, we see negligible effects on prediction performance in
    one output or the other when one output head is essentially turned
    off/down. Because of this, we hold the default value for both prediction
    tasks at 1.0.}
    \label{fig:loss-weight}
\end{figure}

%%%%%% Supplemental Figure 2 %%%%%%
\begin{figure}[H]
    \centering
    \begin{subfigure}[t]{0.49\textwidth}
        \centering
        \includegraphics[width=\textwidth]{figures/suppfig2_LD_lowrecomb.png}
        \caption{}
        \label{fig:ld-decay-low-recomb}
    \end{subfigure}
    \hfill
    \begin{subfigure}[t]{0.49\textwidth}
        \centering
        \includegraphics[width=\textwidth]{figures/suppfig2_LD_highrecomb.png}
        \caption{}
        \label{fig:ld-decay-high-recomb}
    \end{subfigure}
    \caption{\textbf{Linkage disequilibrium under different
    recombination rates for genomic architecture 1
    (\textbf{a)} Recombination = $10^{-8}$;
    \textbf{b)} Recombination = $10^{-6}$).}
    Linkage disequilibrium between selected QTL windows and SNPs outside the
    selected region. To identify SNPs in linkage disequilibrium (LD) with
    the selected QTL window used in genetic architecture 1, we calculated
    pairwise correlations between SNPs inside the selected window and SNPs
    outside the window. These plots do not represent genome-wide LD decay or
    complete LD structure across the simulated genomic element; rather, they
    show the subset of outside-window SNPs that pass the LD threshold with
    SNPs in the selected region. Panel \textbf{a} shows LD partners under lower
    recombination, where LD extends more broadly away from the selected
    window. Under this parameterization, 525 SNP pairs showed very strong LD
    ($r^2 \geq 0.8$), 1,026 showed strong LD ($r^2 \geq 0.5$), and 2,923 showed moderate
    LD ($r^2 \geq 0.2$). Panel \textbf{b} shows the same analysis under higher
    recombination, where LD with the selected window is reduced. Under this
    parameterization, 9 SNP pairs showed very strong LD ($r^2 \geq 0.8$), 24
    showed strong LD ($r^2 \geq 0.5$), and 309 showed moderate LD ($r^2 \geq 0.2$).
    Because SNPs were filtered to include only LD partners of the selected
    QTL window, long-range LD in these panels should be interpreted as
    evidence of SNPs that remain correlated with the selected region, rather
    than as a general description of background LD across the genome.}
    \label{fig:ld-decay}
\end{figure}

%%%%%% Supplemental Figure 3 %%%%%%
\begin{figure}[H]
    \centering
    \includegraphics[width=\textwidth]{figures/suppfig3_permutation.png}
    \caption{\textbf{RMSE and $R^2$ of \ecolocator predictions under
    different permutation regimes (combinations of random \% permuted and
    sites in LD permuted).}
    \ecolocator prediction performance under random and LD-targeted
    permutation regimes ($R^2$ is computed out-of-sample and can be negative
    when predictions are worse than the training-set mean). To assess whether
    prediction accuracy depends
    primarily on the selected QTL window, SNPs in linkage disequilibrium
    with that window, or genome-wide signal distributed across the genotype
    matrix, we permuted subsets of the genotype matrix across individuals
    prior to prediction. The left column shows targeted permutations of the
    selected window alone and the selected window combined with SNPs in
    moderate, high, or all LD with the selected region. The right column
    shows random genome-wide permutations of increasing proportions of SNPs.
    The top row shows root mean squared error (RMSE), and the bottom row
    shows $R^2$ for X coordinate (longitude), Y coordinate (latitude), and the
    environmental covariate. In both permutation designs, prediction
    accuracy declines as more genotype information is shuffled, indicating
    that \ecolocator relies on genetic signal distributed across the genotype
    matrix. However, permuting the selected window together with all
    LD-linked sites causes a sharp loss of performance, suggesting that SNPs
    correlated with the selected region carry important predictive
    information. Randomly permuting all SNPs produces similarly poor
    performance, confirming that prediction accuracy is lost when
    genotype--environment and genotype--location associations are fully
    disrupted. Note that the ``window + all LD'' grouping comprised roughly
    75\% of sites, so comparing that to the random 75\% permutation again
    reiterates the importance of SNPs in LD.}
    \label{fig:permutation}
\end{figure}

%%%%%% Supplemental Figure 4 %%%%%%
\begin{figure}[H]
    \centering
    \includegraphics[width=\textwidth]{figures/suppfig4_localadaptation.png}
    \caption{\textbf{Distribution of strength of local adaptation across parameter space.}
    Kernel density estimates of local adaptation (LA) values for 100
    simulation replicates per parameter combination across three parameter
    explorations on an 8\texttimes 8 grid-like landscape with Hurst = 0 (no
    autocorrelation between space and environment). LA is calculated as the
    mean absolute difference between individual phenotypic values and their
    local optima; lower values indicate stronger local adaptation. The
    selection parameter \(\sigma_K\) controls the standard deviation of the fitness
    function, where smaller values produce a tighter fitness distribution
    and thus stronger selection. Dispersal ($\sigma_D$ = 0.1, 0.3, 0.5; genetic
    architecture 1 and \(\sigma_K\) = 0.5 fixed): increasing dispersal progressively
    eroded local adaptation, with mean LA values of 0.096, 0.130, and 0.146
    for $\sigma_D$ = 0.1, 0.3, and 0.5, respectively. Genetic architecture ($g$ = 1, 2,
    3; $\sigma_D$ = 0.1 and \(\sigma_K\) = 0.5 fixed): distributing QTLs increasingly across the
    genomic element had a modest effect on local adaptation, with mean LA
    values of 0.095, 0.100, and 0.103 for $g$ = 1, 2, and 3, respectively.
    Selection (\(\sigma_K\) = 0.5, 0.7, 0.9; genetic architecture 1 and $\sigma_D$ = 0.3 fixed):
    an extended range of selection strengths was used here to illustrate how 
    local adaptation varies across increasingly weak selection regimes.
    Weaker selection substantially reduced local adaptation, with mean LA
    values of 0.134, 0.187, and 0.228 for \(\sigma_K\) = 0.5, 0.7, and 0.9,
    respectively.}
    \label{fig:local-adaptation}
\end{figure}

%%%%%% Supplemental Figure 5 %%%%%%
\begin{figure}[H]
    \centering
    \includegraphics[width=\textwidth]{figures/suppfig5_downsampling.png}
    \caption{\textbf{Effects of down-sampling individuals and SNPs as input to \ecolocator.}
    \ecolocator performance under various downsampling regimes of the same
    simulated data ($\sigma_D$ = 0.1, \(\sigma_K\) = 0.5, $g$ = 1). In the top row, there are 1000 SNPs
    made available to the model; in the middle and bottom rows, there are
    500 and 100 SNPs respectively. In the columns we vary sample size, with
    100, 500, and 1000 samples from left to right. $R^2$ values across 100
    simulation replicates are summarized in boxplots, with Longitude (x) in
    purple, Latitude (y) in teal, and Environment in yellow. Because
    \ecolocator learns from variation present in the data, it is no surprise
    that we observe the highest prediction accuracy across targets
    (Lon. $R^2$ = 0.954, Lat. $R^2$ = 0.953, and Env. $R^2$ = 0.756) with 1000 samples and
    1000 SNPs (top right corner), and we see the lowest prediction accuracy
    (Lon. $R^2$ = 0.742, Lat. $R^2$ = 0.769, and Env. $R^2$ = 0.493) with only 100
    SNPs and 100 samples (bottom left corner). Under the sample sizes and
    SNP sets explored here, we do not observe diminishing returns with ``too
    much'' data being shown to the model, and similarly, we do not see a
    complete absence of predictive performance when sample size and SNP set
    are restricted.}
    \label{fig:downsampling}
\end{figure}

%%%%%% Supplemental Figure 6 %%%%%%
\begin{figure}[H]
    \centering
    \includegraphics[width=0.5\textwidth]{figures/suppfig6_geoerror.png}
    \caption{\textbf{Distribution of geographic prediction errors on \dfir{} data.}
    A histogram of geographic distance prediction errors in our
    empirical application of \ecolocator on the \dfir case study data.
    Distance in kilometers (km) is on the x-axis and count is on the y-axis. 
    Most errors fall between 0-100km, with others as high as $>$600km.}
    \label{fig:geo-error}
\end{figure}

%%%%%% Supplemental Figure 7 %%%%%%
\begin{figure}[H]
    \centering
    \includegraphics[width=\textwidth]{figures/suppfig7_envmaps.pdf}
    \caption{\textbf{Environmental Maps used in Continuous Space \slim simulations.}
    Environmental maps generated with the \texttt{NLMpy} python package 
    with pixel count set to 8x8 and Hurst (autocorrelation)
    set to H=0. We generated 110 maps and sampled 100 of them randomly
    for each simulation regime. }
    \label{fig:env-maps}
\end{figure}

%%%%%% Supplemental Figure 8 %%%%%%
\begin{figure}[H]
    \centering
    \includegraphics[width=\textwidth]{figures/suppfig8_shap_wza_topcan_aucpr_corrspace.pdf}
    \caption{\textbf{AUC-PR of \ecolocator\texttimes{}SHAP, \texttt{WZA}, and Top-Candidate GEA Methods Across Two Dispersal/Selection Regimes.}
    Precision-recall curves (top) and per-replicate AUC-PR violin plots (bottom) comparing three
    genotype-environment association methods for identifying adaptive SNPs: \ecolocator\texttimes{}SHAP (purple;
    mean $\lvert\mathrm{SHAP}\rvert$ for the environmental covariate, ranked by an empirical FDR
    score computed from known ground-truth adaptive status), \texttt{WZA} (yellow), and top-candidate (blue).
    Top row: precision-recall curves pooling all SNPs across 100 simulation replicates per regime,
    scored by $-\log_{10}$ of each method's $p$-value or empirical FDR score against true adaptive status;
    dashed line = random-baseline AUC-PR (proportion of adaptive SNPs). Bottom row: distribution of
    AUC-PR computed separately within each of the 100 replicates, showing the consistency of each
    method's performance across independent simulations. Left: limited-dispersal, strong-selection
    regime (dispersal=0.1, std.\ dev.\ of fitness function=0.5). Right: higher-dispersal, weaker-selection
    regime (dispersal=0.5, std.\ dev.\ of fitness function=0.7). \ecolocator\texttimes{}SHAP outperforms \texttt{WZA} and
    top-candidate under both regimes, both pooled and at the per-replicate level.}
    \label{fig:gea-comparison-regimes}
\end{figure}

%%%%%% Supplemental Figure 9 %%%%%%
\begin{figure}[H]
    \centering
    \includegraphics[width=\textwidth]{figures/suppfig9_shap_manhattan_dfir_positional.pdf}
    \caption{\textbf{Mean $\lvert\mathrm{SHAP}\rvert$ by Real Genomic Position Across Five Prediction Targets for \dfir.}
    Manhattan-style plots of mean $\lvert\mathrm{SHAP}\rvert$ (averaged across leave-one-out
    cross-validation folds) at real chromosomal positions (chromosomes 1--13, \dfir reference
    assembly DFv5.29; SNP positions resolved by SAM alignment restricted to primary, high-confidence
    mappings, MAPQ $\geq$ 30) for five prediction targets, top to bottom: longitude, latitude, Mean
    Cold Month Temperature (MCMT), Summer Heat:Moisture Index (SHM), and mean summer precipitation
    (MSP). SNPs in the 179 highest-attribution SNPs (approximately the top 1\%) of mean $\lvert\mathrm{SHAP}\rvert$ for a given target are highlighted
    in pink; all other SNPs are shown in blue. Alternating background shading demarcates chromosome
    boundaries.}
    \label{fig:shap-manhattan-dfir}
\end{figure}

%%%%%% Supplemental Figure 10 %%%%%%
\begin{figure}[H]
    \centering
    \includegraphics[width=\textwidth]{figures/suppfig10_LA_greyscale4.png}
    \caption{\textbf{Environmental Map used in Correlated Space Simulations.}
    A grayscale left-right gradient image used in simulations to represent 
    an arbitrary environmental variable correlated with space along the longitudinal axis.}
    \label{fig:env-map-correlated}
\end{figure}

%%%%%% Supplemental Figure 11 %%%%%%
\begin{figure}[H]
    \centering
    \includegraphics[width=\textwidth]{figures/suppfig11_loo_r2_boxplots_both_regimes.png}
    \caption{\textbf{Per-replicate LOOCV prediction accuracy ($R^2$) across two dispersal/selection regimes.}
    For each of 100 independent simulation replicates per regime, 
    predictions from all leave-one-out cross-validation folds 
    were pooled and compared to the known sampled values to 
    compute a single $R^2$ per replicate for each of three prediction 
    targets: latitude, longitude, and environment (cov1). 
    Boxplots show the distribution of these per-replicate $R^2$ 
    values (box = interquartile range, line = median, whiskers = 1.5\texttimes{}IQR, 
    points beyond whiskers = outlier replicates); numbers above each box 
    give the across-replicate mean $R^2$. Left: limited-dispersal, 
    strong-selection regime (\(\sigma_D\)=0.1, \(\sigma_K\)=0.5). Right: higher-dispersal,
    weaker-selection regime (\(\sigma_D\)=0.5, \(\sigma_K\)=0.7). Prediction accuracy
    is uniformly high and tightly distributed for all three targets 
    under limited dispersal/strong selection (mean $R^2$=0.95--0.96), 
    but declines and becomes more variable across replicates under higher 
    dispersal/weaker selection, most notably for latitude 
    (mean $R^2$=0.70 vs. 0.80--0.83 for longitude and environment).}
    \label{fig:loocv-accuracy}
\end{figure}

%%%%%% Supplemental Figure 12 %%%%%%
\begin{figure}[p]
    \centering
    \adjustimage{width=1.3\textwidth, max totalheight=0.9\textheight, center}{figures/suppfig12_shap_manhattan_alltargets_rep1.pdf}
    \caption{\textbf{Genomic Distribution of Mean $\lvert\mathrm{SHAP}\rvert$ Across Three Prediction Targets in a Representative Replicate.}
    Mean $\lvert\mathrm{SHAP}\rvert$ by genomic position (bp) for one representative simulation
    replicate (rep1), shown for three prediction targets (rows: latitude, longitude, environment) and
    both dispersal/selection regimes (columns: limited-dispersal, strong-selection, dispersal=0.1,
    std.\ dev.\ of fitness function=0.5, left; higher-dispersal, weaker-selection, dispersal=0.5,
    std.\ dev.\ of fitness function=0.7, right). SNPs classified as non-neutral (true causal loci; magenta, black-outlined) are overlaid on all other
    neutral/LD-linked SNPs (light blue).}
    \label{fig:shap-manhattan-sim}
\end{figure}
\clearpage

%%%%%% Supplemental Table 1 %%%%%%
\begin{table}[ht]
  \centering
  \caption{\textbf{$R^2$ values for \ecolocator predictions versus Location+ClimateNA estimates.}
  For latitude and longitude, values are \ecolocator's leave-one-out $R^2$, duplicated here for
  reference (ClimateNA does not predict location, hence ``--''). For the three climate variables,
  $R^2$ is shown for \ecolocator's direct climate prediction and for ClimateNA modeled values at
  \ecolocator's predicted locations, computed for the same 348 individuals from a single leave-one-out
  run; $p$-values are from a paired Wilcoxon signed-rank test on absolute error (direct vs.\
  Location+ClimateNA).}
  \label{tab:supptable1}
  \begin{tabular}{@{}p{6cm}ccc@{}}
  \hline
  \textbf{Predicted Location and Climate} & \shortstack{\textbf{Geolocation and}\\ \textbf{Climate Prediction}} & \shortstack{\textbf{Geolocation}\\ \textbf{and ClimateNA}} & \textbf{$p$-value} \\
  \hline
  Latitude (LAT, \textdegree) & 0.83 & -- & -- \\
  Longitude (LON, \textdegree) & 0.75 & -- & -- \\
  Mean Cold Month Temperature (MCMT, \textdegree C) & 0.72 & 0.56 & $1.1\times10^{-6}$ \\
  Mean Summer Precipitation (MSP, mm) & 0.52 & $-0.005$ & $6.0\times10^{-10}$ \\
  Summer Heat:Moisture Index (SHM) & 0.54 & 0.005 & $1.6\times10^{-13}$ \\
  \hline
  \end{tabular}
\end{table}

%%%%%% Supplemental Table 2 %%%%%%
\begin{table}[ht]
  \centering
  \caption{\textbf{Mean $R^2$ values for \ecolocator predictions of longitude, latitude, and environment across simulation types.}}
  \label{tab:supptable2}
  \begin{tabular}{@{}llccc@{}}
  \hline
  \textbf{Simulation type} & \textbf{Parameter} & \textbf{Longitude $R^2$} & \textbf{Latitude $R^2$} & \textbf{Environment $R^2$} \\
  \hline
  Wright--Fisher (discrete space) & -- & 0.986 & 0.985 & 0.916 \\
  \hline
  \multirow{9}{*}{Non-Wright--Fisher (continuous space)}
   & $g=1$ & 0.967 & 0.967 & 0.759 \\
   & $g=2$ & 0.967 & 0.967 & 0.771 \\
   & $g=3$ & 0.968 & 0.968 & 0.787 \\
   \cline{2-5}
   & $\sigma_K=0.3$ & 0.917 & 0.916 & 0.609 \\
   & $\sigma_K=0.5$ & 0.911 & 0.911 & 0.561 \\
   & $\sigma_K=0.7$ & 0.906 & 0.907 & 0.421 \\
   \cline{2-5}
   & $\sigma_D=0.1$ & 0.968 & 0.967 & 0.769 \\
   & $\sigma_D=0.3$ & 0.912 & 0.911 & 0.556 \\
   & $\sigma_D=0.5$ & 0.815 & 0.814 & 0.336 \\
  \hline
  \end{tabular}
\end{table}

%%%%%% Supplemental Table 3 %%%%%%
\begin{table}[ht]
  \centering
  \caption{\textbf{Pairwise overlap among high-attribution SNPs across 
  \ecolocator prediction targets in coastal \dfir.}}
  \label{tab:dfir-shap-overlap}
  \begin{tabular}{@{}lccccc@{}}
  \hline
  \textbf{Prediction target} & \textbf{LAT} & \textbf{LON} & \textbf{MCMT} & \textbf{MSP} & \textbf{SHM} \\
  \hline
  Latitude (LAT)  & 179 & 10 (5.6\%) & 13 (7.3\%) & 33 (18.4\%) & 31 (17.3\%) \\
  Longitude (LON) &     & 179        & 97 (54.2\%) & 5 (2.8\%) & 4 (2.2\%) \\
  MCMT            &     &            & 179         & 5 (2.8\%) & 5 (2.8\%) \\
  MSP             &     &            &             & 179       & 133 (74.3\%) \\
  SHM             &     &            &             &           & 179 \\
  \hline
  \end{tabular}
\end{table}


\end{document}
